\chapter*{专题练习}
\begin{enumerate}
    \item 点集与数集的差别
    \begin{enumerate}
        \item  已知集合$A=\{y | x^2+mx-y+2=0, x \in R\}$， $B=\{y | x+y+1=0, 0 \leq x \leq 1\}$，若$A \cap B \neq \varnothing$，$m$的取值范围是\blkda[5]{$(-\infty, -2\sqrt{3}] \cup [2\sqrt{3}, +\infty)$}. 
        \begin{jieda}
            首先可以判断出来$B = [-2, -1]$，而$A$是函数$y = x^2+mx+2$的函数值的取值范围。由于函数$y = x^2+mx+2$的开口向上，欲使$A \cap B \neq \varnothing$，则应有其函数值最小值小于等于-1。即$\frac{m^2}{4} - \frac{m^2}{2} + 2 \leq -1$，解得$m \in (-\infty, -2\sqrt{3}] \cup [2\sqrt{3}, +\infty)$
        \end{jieda}
        \item  已知集合$A=\{(x, y) | x^2+mx-y+2=0, x \in R\}$， $B=\{(x, y) | x+y+1=0, 0 \leq x \leq 1\}$，若$A \cap B \neq \varnothing$，$m$的取值范围是\blkda[2]{$(-\infty, -5]$}. 
        \begin{jieda}
            根据题意，两集合交集非空代表两段函数的图像有交点，即方程$x^2+(m+1)x+3 = 0$在$[0, 1]$内有根，因为$x = 0$并非原方程的根，故可等价于方程$-(m+1) = x + \frac{3}{x}$在$(0, 1]$内有根 \\
            $y = x + \frac{3}{x}$在$(0, 1]$内单调减，故$-(m+1) \in [4, +\infty)$，故$m \in (-\infty, -5]$
        \end{jieda}
    \end{enumerate}
    \item 退化问题
    \begin{enumerate}
        \item 集合$A = \{x | x^2-8x+12 = 0\}, B = \{x | ax-1 = 0\}$，且$A \cup B = A$，则$a$的集合为\blkda{$\{0, \frac{1}{2}, \frac{1}{6}\}$}
        \begin{jieda}
            $A \cup B = A$即$B \subseteq A$
            \begin{dingautolist}{172}
                \item $a = 0$，此时$B = \varnothing$，满足要求.
                \item $a \neq 0$，此时$B = \{\frac{1}{a}\}$，因为$A = \{2, 6\}$，若$B \subseteq A$，则$\frac{1}{a} \in A$，因此$a = \frac{1}{2} \mbox{或} \frac{1}{6}$
            \end{dingautolist}
        \end{jieda}
        \item 对任意实数$x$，不等式$2kx^2+kx-3<0$恒成立，则实数$k$的取值范围是\blkda{$(-24, 0]$}
        \begin{jieda}
            当$k=0$时，不等式成立。\\
            当$k \neq 0$，从二次函数的角度看一元二次不等式，要使得不等式恒成立，则应有函数开口向下，即$k<0$，并且有$\Delta = k^2 + 24k < 0$，因此$k \in (-24, 0)$。\\
            综上，$k \in (-24, 0]$
        \end{jieda}
        \item 已知集合$A = \{(x, y) | \frac{y-3}{x-2} = a+1\}, B = \{(x, y) | (a^2-1)x+(a-1)y = 15\}$，若$A \cap B = \varnothing$，则的取值范围是\blkda[3]{$\{1, -1, -4, \frac{5}{2}\}$}
        \begin{jieda}
            $A = \{(x, y) | y = 3 + (a+1)(x-2), x \neq 2 \} \\
            B = \{(x, y) | y = -(a+1)x + \frac{15}{a-1}\} (a \neq 1), \quad B = \varnothing(a = 1)$ \\
            易知当$a = 1$时符合条件 \\
            当$a \neq 1$时，由于$A$是一条直线刨除了点$(2,3)$，$B$是另一条直线，所以有两种可能满足题意。
            \begin{enumerate}
                \item 两直线平行，此时$a = -1$，满足题意
                \item $B$对应直线过$(2, 3)$，此时$a = -4 \mbox{或} \frac{5}{2}$，均满足题意
            \end{enumerate}
        \end{jieda}
    \end{enumerate}
    \item 包含关系与推出关系
    \begin{enumerate}
        \item 已知集合$A = \{x | \frac{x-2}{x+1} < 0\}$，$x \in A$的一个必要条件是$x \geq a$，则实数$a$的取值范围是\blkda{$a \leq -1$}
        \begin{jieda}
            $A = (-1, 2)$，故$a \leq -1$ \\
            \ps \textbf{小范围可以推出大范围}
        \end{jieda}
        \item 已知$x \in \mathbf{R}$，“$(x-2)(x-3) \leq 0$”是“$|x-2|+|x-3| = 1$成立”的\blkda{充要}条件.
        \begin{jieda}
            $(x-2)(x-3) \leq 0$解集为$[2, 3]$ \\
            $|x-2|+|x-3| = 1$解集为$[2, 3]$（平底锅函数，三角不等式）
        \end{jieda}
    \end{enumerate}
    \item 求解分式不等式
    \begin{enumerate}
        \item 不等式$\frac{x+5}{x^2+2x+3} \geq 1$的解集为\blkda{$[-2, 1]$}
        \begin{jieda}
            因为$x^2+2x+3 > 0$恒成立，因此可以转化为$x+5 \geq x^2+2x+3$，解集为$[-2, 1]$ \\
            \ps 先判断分母是否恒正或者恒负，再考虑移项通分.
        \end{jieda}
        \item 不等式$\frac{x^2(x+3)(x^2-x+1)}{x^2 - 5x + 4} > 0$的解集为\blkda[5]{$(-3, 0) \cup (0, 1) \cup (4, +\infty)$}
        \begin{jieda}
            \ps 穿针引线法求解，奇穿偶回，\textbf{偶次方根绝对不能丢弃，且易出错}.
        \end{jieda}
        \item 定义区间$(c, d), [c, d), (c, d], [c, d]$的长度均为$d-c (d>c)$。已知实数$a > b$，则满足$\frac{1}{x-a} + \frac{2}{x-b} \geq 3$的$x$构成的区间长度之和为\blkda{1}
        \begin{jieda}
            首先通分有$\frac{3x-2a-b}{(x-a)(x-b)} \geq 3$，易知移项通分运算量很大。因此考虑分类讨论乘分母
            \begin{enumerate}
                \item $x > a$ 或 $x < b$ \\
                此时不等式可以转化为$3x-2a-b \geq 3(x-b)(x-a)$
                \item $a > x > b$ \\
                此时不等式可以转化为$3x-2a-b \leq 3(x-b)(x-a)$
            \end{enumerate}
            绘制$y = 3(x-b)(x-a)$和$y = 3x-2a-b$的草图可知不等式的解集为$(b, x_1) \cup (a, x_2)$，其中$x_1, x_2$是方程$3(x-b)(x-a) = 3x-2a-b$的两根。根据韦达定理可知解集的区间长度和为$x_1 + x_2 - a - b = 1$ \\
            \ps 本题也可以通过绘制$y = \frac{1}{x-a} + \frac{2}{x-b}$的草图，分析其单调区间，值域来解决。本质是一样的。
        \end{jieda}
    \end{enumerate}
    \item 求解含绝对值不等式
    \begin{enumerate}
        \item 解不等式：$|x+1| < 3-2x$
        \begin{jieda}
            解集为$(-\infty, \frac{2}{3})$ \\
            \ps \begin{enumerate*}[1.]
                \item $|f(x)| > g(x)$的充要条件是$f(x) > g(x)$或$f(x) < -g(x)$ 
                \item $|f(x)| < g(x)$的充要条件是$-g(x) < f(x) < g(x)$
            \end{enumerate*} \\
            \ps \textbf{此处绝对不能平方处理}，否则等价于$|x+1| < |3-2x|$会得到$(-\infty, \frac{2}{3}) \cup (4, +\infty)$；当不等式两侧均含有绝对值时才适合平方去绝对值.
        \end{jieda}
        \item 若关于$x$的不等式$ax + 6 + |x^2 - ax - 6| \geq 4$恒成立，$a$的取值范围是\blkda{$[-1, 1]$}
        \begin{jieda}
            原不等式转化为$|x^2 - ax - 6| \geq -2 - ax$恒成立，即$x^2 - ax - 6 \geq -2 - ax$与$x^2 - ax - 6 \leq 2 + ax$解集的并集为$\mathbf{R}$ \\
            前者解集为$(-\infty, -2] \cup [2, +\infty)$，后者解集为$[a-\sqrt{a^2+8}, a + \sqrt{a^2+8}]$，故$(-2, 2) \subseteq [a-\sqrt{a^2+8}, a + \sqrt{a^2+8}]$，因此$a \in [-1, 1]$
        \end{jieda}
        \item 关于实数$x$的不等式$\left | x - \frac{(a+1)^2}{2} \right | \leq \frac{(a-1)^2}{2}$与$|x-a-1| \leq a$的解集依次是$A, B$，求使得$A \subseteq B$的$a$的取值范围.
        \begin{jieda}
            $\left | x - \frac{(a+1)^2}{2} \right | \leq \frac{(a-1)^2}{2}$即$-\frac{(a-1)^2}{2} \leq x - \frac{(a+1)^2}{2} \leq \frac{(a-1)^2}{2}$，即$2a \leq x \leq a^2+1$ \\
            $|x-a-1| \leq a$ 即 $-a \leq x-a-1 \leq a$，即$1 \leq x \leq 2a+1$ \\
            由于$a^2+1 \geq 2a$恒成立，所以$A$必不为空集，故$A \subseteq B$即$1 \leq 2a, a^2+1 \leq 2a+1$，解得$a \in [\frac{1}{2}, 2]$
        \end{jieda}
        \item 已知函数$f(x) = |x-2| + |x-a|$，若对于任意的$x \in [1, 2]$，$f(x) \geq |x-4|$恒成立，实数$a$的范围是\blkda[4]{$(-\infty, -1] \cup [4, +\infty)$}
        \begin{jieda}
            当$x \in [1, 2]$，$f(x) \geq |x-4|$即$2-x + |x-a| \geq 4-x$，整理得$|x-a| \geq 2$，由绝对值的几何意义可知$a \in (-\infty, -1] \cup [4, +\infty)$ \\
            \ps 有自变量的具体范围的时候，可以先看一下绝对值是否可以直接去掉.
        \end{jieda}
        \end{enumerate}
    \item 整根问题
    \begin{enumerate}
        \item 若关于$x$的不等式$k|x| > |x - 2|$恰有4个整数解，则实数$k$的取值范围是\blkda{$(\frac{3}{5}, \frac{2}{3}]$}
        \begin{jieda}
            易知$x = 0$不是不等式的解，故$k|x| > |x - 2|$进行参变分离，同解变形$k > |\frac{x-2}{x}|$，即$k > |1 - \frac{2}{x}|$ \\
            由函数$f(x) = |1 - \frac{2}{x}|$的图像可知整根的集合为$\{2, 3, 4, 5\}$，故$f(5) < k \leq f(6)$，即$k \in (\frac{3}{5}, \frac{2}{3}]$
            \ps \textbf{解决整根问题的第一步往往是确认整根的集合}
        \end{jieda}
        \item 已知关于$x$的不等式组$\left \{ \begin{array}{l}
             \frac{x-a}{x-1+a} < 0\\
             x+3a>1 
        \end{array}\right.$的整数解恰有两个，则实数$a$的取值范围是\blkda{$(1, 2]$}
        \begin{jieda}
            首先可以确定第一个不等式的解集为$(a, 1-a)$或$(1-a, a)$，因为$a$和$1-a$关于$\frac{1}{2}$对称，因此$\frac{1}{2}$必然在该不等式的解集中。而第二个不等式的解集为$(1-3a, +\infty)$
            \begin{dingautolist}{172}
                \item 若$a \in [0, 1]$，则$(a, 1-a)$或$(1-a, a)$没有整数解，故整个不等式组的解集中也没有整数解.
                \item 若$a < 0$，则$a < 0 < 1 < 1-a < 1-3a$，则整个不等式组解集为空集.
                \item 若$a > 1$，则$1-3a < 1-a < 0 < 1 < a$整个不等式组的解集为$(1-a, a)$，此时解集中必有两个整数解$0, 1$，故当$a \in (1, 2]$时，满足要求.
            \end{dingautolist}
            综上，$a \in (1, 2]$ \\
            \ps 本题画数轴可以更方便分析，本题判断出$a$和$1-a$关于$\frac{1}{2}$对称后可以猜测到两个整根是0和1.
        \end{jieda}
        \item 已知不等式组$ \left\{ \begin{array}{l}
             x^2-x+a-a^2 < 0  \\
             x + 2a > 1
        \end{array} \right.$的整数解恰好有两个，$a$的取值范围是\blkda{$(1, 2]$}
        \begin{jieda}
            设函数$f(x) = x^2 - x$，则第一个不等式即$f(x) < f(a)$，根据函数的性质可知不等式即$|x-\frac{1}{2}| < |a - \frac{1}{2}|$，故$\frac{1}{2}-|a - \frac{1}{2}| < x < \frac{1}{2} + |a - \frac{1}{2}|$，故第一个不等式解集为$(1 - a, a)$或$(a, 1-a)$，第二个不等式解集为$(1-2a, +\infty)$
            \begin{dingautolist}{172}
                \item $a \in [0, 1]$，此时第一个不等式中不含有两个整数解，不符合要求.
                \item $a < 0$，第一个不等式解集为$(a, 1-a)$，而$1-a < 1-2a$，故不等式组解集为空集，不符合要求.
                \item $a > 1$，第一个不等式解集为$(1 - a, a)$，故不等式解集为$(1 - a, a)$，当$a \in (1, 2]$满足要求.
            \end{dingautolist}
        \end{jieda}
    \end{enumerate}
    \item 不等式恒成立问题
    \begin{enumerate}
        \item 关于$x$的不等式$x^2+9+|x^2-3x| \geq kx$在区间$[1,5]$上恒成立，则实数$k$的取值范围是\blkda{$k \leq 6$}
        \begin{jieda}
            当$x \in [1, 5]$时，原不等式可以转化为$x + \frac{9}{x} + |x - 3| \geq k$ \\
            不等式恒成立即$k$比左侧最小值小，而易知左侧最小值在$x = 3$时取，故$k \leq 6$
        \end{jieda}
        \item 若不等式$x^2+px>4x+p-3$，当$0 \leq p \leq 4$时恒成立，则$x$的取值范围是\blkda[4]{$(-\infty, -1) \cup (3, +\infty)$}
        \begin{jieda}
            \begin{enumerate}
                \item 方法1：看做是关于$p$的一元一次不等式求解 \\
                不等式移项得到$x^2 -4x + 3 > (1-x)p$。\\
                因式分解有$(x-1)(x-3) > -p(x-1)$ \\
                下面根据$x$的取值范围进行分类讨论：
                \begin{dingautolist}{172}
                    \item $x > 1$，不等式转化为$x - 3 > -p$，故$x \in (3, +\infty)$
                    \item $x = 1$，这种情况不成立
                    \item $x < 1$，不等式转化为$x - 3 < -p$，故$x \in (-\infty, -1)$
                \end{dingautolist}
                综上，当$x \in (-\infty, -1) \cup (3, +\infty)$时不等式恒成立。
                \item 方法2：看做是关于$x$的一元二次不等式求解 \\
                不等式移项得到$x^2 + (p-4)x + (3-p) > 0$。即$(x-(3-p))(x-1) > 0$ \\
                因此根据两解可以写出不等式的解集为$\left \{
                    \begin{array}{ll}
                         (-\infty, 3-p) \cup (1, +\infty) & 4 \geq p \geq 2 \\
                         (-\infty, 1) \cup (3-p, +\infty) & 0 \leq p < 2
                    \end{array}
                \right. $ \\
                要使得不等式在$p$取到$[0, 4]$内任意值时成立，则应取所有$p$所对应的解集的交集，故当$x \in (-\infty, -1) \cup (3, +\infty)$时不等式恒成立。
            \end{enumerate}
            \ps 通过对比可知以$p$为主元更容易理解和推广。
        \end{jieda}
        \item 已知不等式$x^2+(a-2)x+3a>0(*)$
        \begin{enumerate}
            \item 若不等式$(*)$对任意的$x \in [-1, 1]$恒成立，求实数$a$的取值范围
            \item 若不等式$(*)$对任意的$a \in [-1, 1]$恒成立，求实数$x$的取值范围
        \end{enumerate}
        \begin{jieda}
            \begin{enumerate}
                \item 
                \begin{dingautolist}{172}
                    \item 通解法 \\
                    $f(x) = x^2+(a-2)x+3a$的对称轴为$x = 1-\frac{a}{2}$ \\
                    若$1-\frac{a}{2} < -1$，即$a > 4$，此时$f(x)$在$[-1, 1]$上单调递增，因此$f(-1) > 0$，即$1+2-a+3a > 0$，解得$a > -\frac{3}{2}$，结合对称轴有$a > 4$ \\
                    若$1-\frac{a}{2} > 1$，即$a < 0$，此时$f(x)$在$[-1, 1]$上单调递减，因此$f(1) > 0$，即$1+a-2+3a > 0$，解得$a > \frac{1}{4}$，此种情况无解 \\
                    若$1-\frac{a}{2} \in [-1, 1]$，即$a \in [0, 4]$，此时$f(x)$在$[-1, 1]$内最小值在顶点处取，即$f(1-\frac{a}{2}) > 0$，整理有$a^2-16a+4<0$，解得$a \in (8-2\sqrt{15}, 8+2\sqrt{15})$，结合对称轴有$a \in (8-2\sqrt{15}, 4]$ \\
                    综上，$a \in (8-2\sqrt{15}, +\infty)$
                    \item 参变分离 \\
                    $(x+3)a-2x+x^2 > 0$在$x \in [-1, 1]$恒成立，即$a > \frac{2x-x^2}{x+3}$在$x \in [-1, 1]$恒成立.\\
                    换元$t = x+3$， $t \in [2, 4]$，即$a > \frac{-t^2+8t-15}{t}$在$[2, 4]$恒成立.进一步变形有$a > -(t+\frac{15}{t})+8$在$[2, 4]$恒成立.根据对勾函数性质可知$-(t+\frac{15}{t})+8 \leq 8 - 2\sqrt{15}$（$t = \sqrt{15}$，即$x = \sqrt{15}-3$时取等号）,故$a \in (8-2\sqrt{15}, +\infty)$
                \end{dingautolist}
                \item 将$a$做为主元，即$(x+3)a-2x+x^2 > 0$在$a \in [-1, 1]$恒成立，设$f(a) = (x+3)a-2x+x^2$，故只需$f(1) > 0, f(-1) > 0$即可.解得$x \in (-\infty, \frac{3-\sqrt{21}}{2})\cup(\frac{3+\sqrt{21}}{2}, +\infty)$
            \end{enumerate}
            \ps 恒成立问题，首先考虑主元应该是谁，然后考虑是否要参变分离.
            \end{jieda}
    \end{enumerate}
    \item 不等式有解问题
    \begin{enumerate}
        \item 关于$x$的不等式$x^2+9+|x^2-3x| \geq kx$在区间$[1,5]$上有解，则实数$k$的取值范围是\blkda{$k \leq 12$}
        \begin{jieda}
            当$x \in [1, 5]$时，原不等式可以转化为$x + \frac{9}{x} + |x - 3| \geq k$ \\
            不等式有解即$k$比左侧最大值小，设$f(x) = x + \frac{9}{x} + |x - 3|$易知$f(x)$在$[1, 5]$上先减后增，最大值在端点处取. $f(1) = 12, f(5) = \frac{44}{5}$故$k \leq 12$\\
            \ps \textbf{区分有解问题和恒成立问题}.
        \end{jieda}
        \item 已知函数$f(x) = |x + \frac{1}{x} - 4|$，若关于$x$的不等式$f(x) \geq m^2-m+2$在区间$[\frac{1}{6}, 3]$总有解，则实数$m$的取值范围是\blkda[3.5]{$[\frac{3-\sqrt{15}}{6}, \frac{3+\sqrt{15}}{6}]$}
        \begin{jieda}
            首先要意识到不等式的右侧没有$x$，因此可以将右侧当做一个整体，直接找$f(x)$在区间$[\frac{1}{6}, 3]$上的最大值. 根据对勾函数的性质可知$f(x)$在$[\frac{1}{6}, 1]$上先减后增，在$[1, 3]$单调递减. 故最大值会在$x = \frac{1}{6}$或$x = 1$处取到.$f(\frac{1}{6}) = \frac{13}{6}, f(1) = 2$，故$f(x)$的最大值为$\frac{13}{6}$ \\
            问题转化为求解不等式$\frac{13}{6} \geq m^2-m+2$，解集为$[\frac{3-\sqrt{15}}{6}, \frac{3+\sqrt{15}}{6}]$. 
        \end{jieda}
        
    \end{enumerate}
    \item 二次方程根的分布
    \begin{enumerate}
        \item 已知$a, b$为常数，函数$f(x) = x^2 - bx + a$，当$a = 2b-1$时，若方程$f(x) = 0$在$(-2, 1)$上有解，求实数$b$的取值范围.
        \begin{jieda}
            $f(x) = x^2 - bx + 2b - 1$.
            \begin{dingautolist}{172}
                \item 分对称轴位置讨论 \\
                $f(x)$的对称轴为$x = \frac{b}{2}$，$f(-2) = 4b+3$，$f(1) = b$，$\Delta = b^2 - 4(2b-1) = b^2 - 8b + 4$
                \begin{enumerate}[1.]
                    \item 若$\frac{b}{2} \leq -2$，则$f(x)$在$(-2, 1)$单调增，故“方程$f(x) = 0$在$(-2, 1)$上有解”即$f(-2) < 0, f(1) > 0$，此种情况解集为空集.
                    \item 若$\frac{b}{2} \geq 1$，则$f(x)$在$(-2, 1)$单调减，故“方程$f(x) = 0$在$(-2, 1)$上有解”即$f(-2) > 0, f(1) < 0$，此种情况解集为空集.
                    \item 若$\frac{b}{2} \in (-2, 1)$，则$f(x)$在$(-2, 1)$先减后增，故“方程$f(x) = 0$在$(-2, 1)$上有解”即$\Delta \geq 0$且有$f(-2) > 0$或$f(1) > 0$，此种情况解集为$(-\frac{3}{4}, 4-2\sqrt{3}]$
                \end{enumerate}
                综上，$b \in (-\frac{3}{4}, 4-2\sqrt{3}]$
                \item 参变分离 \\
                “方程$f(x) = 0$在$(-2, 1)$上有解”即“$b = \frac{1-x^2}{2-x}$在$(-2, 1)$上有解”，即“$b = \frac{3}{x-2} + (x-2) + 4$在$(-2, 1)$上有解” \\
                因此$b$的取值范围即$y = \frac{3}{x-2} + (x-2) + 4$在$(-2, 1)$上的值域. \\
                $y = \frac{3}{x-2} + (x-2) + 4$在$(-2, \sqrt{3}-2)$单调递增，在$[2-\sqrt{3}, 1)$单调递减，值域为$(-\frac{3}{4}, 4-2\sqrt{3}]$，故$b \in (-\frac{3}{4}, 4-2\sqrt{3}]$
            \end{dingautolist}
            \ps 有解问题可以采用分对称轴位置讨论或参变分离解决，\textbf{小题推荐参变分离，大题可以分类讨论写步骤，参变分离检验}。分对称轴位置讨论时，对于对称轴在目标区间内的情况要注意其充要条件。参变分离会把方程有解问题转化成函数求值域问题。
        \end{jieda}
        \item 已知函数$f(x) = x^2 + 2mx + 2m + 3$至少有一个零点在区间$(0, 2)$内，求实数$m$的取值范围.
        \begin{jieda}
            \begin{dingautolist}{172}
                \item 分对称轴位置讨论 \\
                $f(x)$对称轴为$x = -m$，$f(0) = 2m+3$，$f(2) = 7+6m$，$\Delta = 4m^2 - 8m - 12$
                \begin{enumerate}[1.]
                    \item 若$-m \leq 0$，则$f(x)$在$(0, 2)$单调增，故“$f(x)$在$(0, 2)$上有零点”即$f(0) < 0, f(2) > 0$，此种情况无解.
                    \item $-m \geq 2$，则$f(x)$在$(0, 2)$单调减，故“$f(x)$在$(0, 2)$上有零点”即$f(0) > 0, f(2) < 0$，此种情况无解.
                    \item $-m \in (0, 2)$，则$f(x)$在$(0, 2)$先减后增，“$f(x)$在$(0, 2)$上有零点”即$\Delta \geq 0$，且有$f(0) > 0$或$f(2) > 0$，此种情况解得$m \in (-\frac{3}{2}, -1]$
                \end{enumerate}
                综上，$m \in (-\frac{3}{2}, -1]$
                \item 参变分离
                $f(x)$至少有一个零点在区间$(0, 2)$内，即$f(x) = 0$在$(0, 2)$内有根，即方程$m = \frac{-3-x^2}{2x+2}$在$(0, 2)$内有根，即方程$m = \frac{1}{2} \cdot \frac{-3-(t-1)^2}{t}$在$t \in (1, 3)$内有根，即方程$m = \frac{1}{2} \cdot [-(t+\frac{4}{t})+2]$在$t \in (1, 3)$内有根.\\
                因此$m$的取值范围即$y = \frac{1}{2} \cdot [-(t+\frac{4}{t})+2]$在$t \in (1, 3)$内的值域. \\
                $y = \frac{1}{2} \cdot [-(t+\frac{4}{t})+2]$在$(1, 2)$单调增，在$(2, 3)$单调减，因此值域为$(-\frac{3}{2}, -1]$，故$m \in (-\frac{3}{2}, -1]$
            \end{dingautolist}
        \end{jieda}
        \item 已知函数$f(x) = x^2 + 2mx + 2m + 3$恰有一个零点在区间$(0, 2)$内，求实数$m$的取值范围.
        \begin{jieda}
            \begin{dingautolist}{172}
                \item 通解法 \\
                $f(x)$对称轴为$x = -m$，$f(0) = 2m+3$，$f(2) = 7+6m$，$\Delta = 4m^2 - 8m - 12$
                \begin{enumerate}[1.]
                    \item 符合函数零点存在定理，此时有$f(0) \cdot f(2) < 0$，即$m \in (-\frac{3}{2}, -\frac{7}{6})$
                    \item 二次函数与$x$轴相切在区间内，此时有$\Delta = 0$，且$-m \in (0, 2)$即$m = -1$
                    \item 左端点函数值为0，此时有$m = -\frac{3}{2}$，此时$f(2) < 0$，$f(x)$在$(0, 2)$内先减后增，故不存在零点.
                    \item 右端点函数值为0，此时有$m = -\frac{7}{6}$，此时$f(0) > 0$，$f(x)$在$(0, 2)$内先减后增，存在唯一零点.
                \end{enumerate}
                综上，$m \in (-\frac{3}{2}, -\frac{7}{6}] \cup \{-1\}$
                \item 参变分离
                $f(x)$有一个零点在区间$(0, 2)$内，即方程$m = \frac{-3-x^2}{2x+2}$在$(0, 2)$内有一根，即方程$m = \frac{1}{2} \cdot \frac{-3-(t-1)^2}{t}$在$t \in (1, 3)$内有一根，即方程$m = \frac{1}{2} \cdot [-(t+\frac{4}{t})+2]$在$t \in (1, 3)$内有一根，即函数$y = m$与函数$y = \frac{1}{2} \cdot [-(t+\frac{4}{t})+2]$在$t \in (1, 3)$内只有一个交点.\\
                $y = \frac{1}{2} \cdot [-(t+\frac{4}{t})+2]$在$(1, 2)$单调增，在$(2, 3)$单调减，因此当$m \in (-\frac{3}{2}, -\frac{7}{6}] \cup \{-1\}$时满足要求.
            \end{dingautolist}
            \ps 二次函数在区间内只有一个零点的通解法对应着\textbf{四种情况}：\begin{enumerate*}[1.]
                \item 符合函数零点存在定理
                \item $\Delta = 0$且对称轴在区间内
                \item 左端点为函数的一个零点（应根据区间开闭情况单独检验）
                \item 右端点为函数的一个零点（应根据区间开闭情况单独检验）
            \end{enumerate*} \\
            参变分离会将问题转化为动直线和定曲线只有唯一交点，要\textbf{注意相切的情况和端点情况}。\\
            这种情况参变分离和通解法都没有明显优势.
        \end{jieda}
        \item 已知函数$f(x) = x^2 + 2mx + 2m + 3$恰有一个零点在区间$(0, +\infty)$内，求实数$m$的取值范围.
        \begin{jieda}
            \begin{dingautolist}{172}
                \item 通解法 \\
                $f(x)$对称轴为$x = -m$，$f(0) = 2m+3$，$\Delta = 4m^2 - 8m - 12$
                \begin{enumerate}[1.]
                    \item “符合函数零点存在定理”，此时有$f(0) < 0$，即$m \in (-\infty, -\frac{3}{2})$
                    \item 二次函数与$x$轴相切在区间内，此时有$\Delta = 0$，且$-m \in (0, +\infty)$即$m = -1$
                    \item 左端点函数值为0，此时有$m = -\frac{3}{2}$，$f(x)$在$(0, +\infty)$内先减后增，存在唯一零点.
                \end{enumerate}
                综上，$m \in (-\infty, -\frac{3}{2}] \cup \{-1\}$
                \item 参变分离
                $f(x)$有一个零点在区间$(0, +\infty)$内，即方程$m = \frac{-3-x^2}{2x+2}$在$(0, +\infty)$内有一根，即方程$m = \frac{1}{2} \cdot \frac{-3-(t-1)^2}{t}$在$t \in (1, +\infty)$内有一根，即方程$m = \frac{1}{2} \cdot [-(t+\frac{4}{t})+2]$在$t \in (1, +\infty)$内有一根，即函数$y = m$与函数$y = \frac{1}{2} \cdot [-(t+\frac{4}{t})+2]$在$t \in (1, +\infty)$内只有一个交点.\\
                $y = \frac{1}{2} \cdot [-(t+\frac{4}{t})+2]$在$(1, 2)$单调增，在$(2, +\infty)$单调减，因此当$m \in (-\infty, -\frac{3}{2}] \cup \{-1\}$时满足要求.
            \end{dingautolist}
            \ps 区间一侧为无穷的情况对于通解法可以根据开口方向将问题简化.
        \end{jieda}
        \item 已知$g(x) = -x^2+2x$，若关于$x$的方程$g(x) = -mx-m$在$(0, 2)$上恰有两个不等实根，求$m$的取值范围.
        \begin{jieda}
            \begin{dingautolist}{172}
                \item 通解法 \\
                “方程$g(x) = -mx-m$在$(0, 2)$上恰有两个不等实根”即“方程$-x^2+(2+m)x+m = 0$在$(0, 2)$上恰有两个不等实根”，设$h(x) = -x^2+(2+m)x+m$，对称轴为$x = \frac{2+m}{2}$，$\Delta = (2+m)^2 + 4m$，$h(0) = m$，$h(2) = 3m$ \\
                $h(x) = 0$在$(0, 2)$上恰有两个不等实根的充要条件是$\left \{\begin{array}{l}
                    \Delta > 0  \\
                    \frac{2+m}{2} \in (0, 2) \\
                    h(0) < 0 \\
                    h(2) < 0
                \end{array} \right.$，解得$m \in (2\sqrt{3}-4, 0)$
                \item 参变分离 \\
                “方程$g(x) = -mx-m$在$(0, 2)$上恰有两个不等实根”即“方程$m = \frac{x^2-2x}{x+1}$在$(0, 2)$上恰有两个不等实根”，即“函数$y = m$与函数$y = t + \frac{3}{t} - 4$在$(1, 3)$上恰有两个交点” \\
                $y = t + \frac{3}{t} - 4$在$(1, \sqrt{3})$单调减，在$[\sqrt{3}, 3)$上单调增，因此当$m \in (2\sqrt{3}-4, 0)$时满足要求.
            \end{dingautolist}
        \end{jieda}
    \end{enumerate}
    \item 三角不等式
    \begin{enumerate}
        \item ``$ab\geq 0$''是``$\abs{a+b}=\abs{a}+\abs{b}$''的\blk 条件. \dotfill(\kh{C})
        \xx
        {充分不必要条件}
        {必要不充分条件}
        {充要条件}
        {既不充分也不必要条件}
        \begin{jieda}
        考察三角不等式的\textbf{取等条件}. \\
        $\abs{a+b} \leq \abs{a}+\abs{b}$的取等条件是$ab\geq 0$ \\
        $\abs{a-b} \leq \abs{a}+\abs{b}$的取等条件是$ab\leq 0$
        \end{jieda}
        \item 已知函数$f(x) = |x+m| + |x-2m|$，若$f(x) \geq 3$恒成立，求$m$的取值范围
        \begin{jieda}
            问题可以转化为$f(x)$的最小值大于等于3恒成立，而 $f(x)$是一个平底锅函数，最小值可以通过三角不等式求得，即 $f(x) = |x+m| + |x-2m| \geq |3m|$，因此有$|3m| \geq 3$，得到$m \in (-\infty, -1] \cup [1, +\infty)$ \\
            \ps 注意\textbf{问题转化}的逻辑；注意两个绝对值的式子相加减，想一下三角不等式
        \end{jieda}
        \item 若实数$a$使得不等式$|x-2a| + |2x-a| \geq a^2$对任意的$x$恒成立，实数$a$的取值范围是\blkda{$[-\frac{3}{2}, \frac{3}{2}]$}
        \begin{jieda}
            \begin{enumerate}
            \item 方法1 \\
            首先将问题转化为左侧的最小值大于等于右侧，设函数$f(x) = |x-2a| + |2x-a|$，根据斜底锅函数的图象可知，其最低点必然位于系数绝对值最大的一项取零的位置，即$x = \frac{a}{2}$，最小值为$|\frac{3a}{2}|$，由此不等式可以转化为$|\frac{3a}{2}| \geq a^2$，即$|\frac{3}{2}| \geq |a|$，故$a \in [-\frac{3}{2}, \frac{3}{2}]$
            \item 方法2 \\
            通过设$x = ka(k \in \mathbf{R})$来简化计算。\\
            根据$x = ka(k \in \mathbf{R})$，不等式可以化为$|ka - 2a| + |2ka - a| \geq a^2$，\\
            消去$|a|$，不等式变为$|k - 2| + |2k - 1| \geq |a|$。\\
            通过对$k$进行分类讨论，可以得到$|k - 2| + |2k - 1|$的取值范围为$[\frac{3}{2}, +\infty)$，故$|a| \leq \frac{3}{2}$，即$a \in [-\frac{3}{2}, \frac{3}{2}]$
            \end{enumerate}
            \ps 斜底锅函数$f(x) = |ax+b|+|cx+d|$（$|a| \neq |c|$），若$|a| > |c|$，则取最小值时$x = -\frac{b}{a}$，若$|a| < |c|$，则取最小值时$x = -\frac{d}{c}$
        \end{jieda}
    \item 
    \begin{enumerate}
        \item  若关于$x$的不等式$\abs{x-1}-\abs{x-2}<a$的解集为$\bR$, 则实数$a$的取值范围为\blkda{$a>1$}
        
        \item 若关于$x$的不等式$\abs{x-1}-\abs{x-2}<a$有实数解, 则实数$a$的取值范围为\blkda{$a>-1$}. 
        
        \item 若关于$x$的不等式$\abs{x-1}-\abs{x-2}<a$的解集为$\varnothing$, 则实数$a$的取值范围为\blkda{$a\leq-1$}. 
    \end{enumerate}
    \begin{jieda}
        斜坡函数$f(x) = |x+a|-|x+b|$，最大值为$|a-b|$，最小值为$-|a-b|$. \\
        例如：$f(x) = |x+1|-|x+2| = \left \{ \begin{array}{ll}
            1 & x < -2 \\
            -2x-3 & x \in [-2, -1] \\
            -1 & x > -1 \\
        \end{array}\right.$
    \end{jieda}
    \end{enumerate}
    \item 无限制条件利用平均值不等式求最值
    \begin{enumerate}
        \item 下列各式中, 最小值为4的是\dotfill(\kh{B})
        \xx
        {$x+2\sqrt{x}+5$}
        {$\frac{x^2+5}{\sqrt{x^2+1}}$}
        {$x^2+3+\frac{4}{x^2+3}$}
        {$x+\frac{4}{x}$}
        \begin{jieda}
            注意取等条件是否可以取到！
        \end{jieda}
        \item 若$a>b>0$, 则$a^2+\frac{16}{b(a-b)}$的最小值是\blkda{16}. 
        \begin{jieda}
        先考虑消元，$b(a-b) \leq \left (\frac{b+(a-b)}{2}\right )^2$，取等条件为$a = 2b$. \\
        故$a^2+\frac{16}{b(a-b)} \geq a^2 + \frac{64}{a^2} \geq 16$. 取等条件为$a = 2\sqrt{2}, b = \sqrt{2}$
        \end{jieda}
        \item 已知不等式$(x+y)\left(\frac{a}{x}+\frac{1}{y}\right)\geq 9$对任意正实数$x,y$恒成立, 则正实数$a$的最小值是\dotfill(\kh{B})
        \xx
        {2}
        {4}
        {6}
        {8}
        \begin{jieda}
            首先求出$(x+y)\left(\frac{a}{x}+\frac{1}{y}\right)$最小值(含$a$)，然后令其大于等于9，得到$a$的取值范围，从中选取最小值\\
            $(x+y)\left(\frac{a}{x}+\frac{1}{y}\right) = a + 1 + \frac{ay}{x} + \frac{x}{y} \geq a + 1 + 2\sqrt{a}$（$x  = \sqrt{a}y$时取等）\\
            设$f(a) = a+1+2\sqrt{a}$，显然$f(a)$单调递增，且有$f(4) = 9$. 故$a \geq 4$时原始$\geq 9$，因此$a$最小值为4.
        \end{jieda}
        \item 已知两两不同的$x_1, y_1, x_2, y_2, x_3, y_3$满足$x_1 + y_1 = x_2 + y_2 = x_3 + y_3$，且$x_1 < y_1, \quad x_2 < y_2, \quad x_3 < y_3, \quad x_1y_1+x_3y_3=2x_2y_2 > 0$，则下列选项中恒成立的是 \dotfill(\kh{B})
        \xx{$2x_2 < x_1 + x_3$}{$2x_2 > x_1 + x_3$}{$x_2^2 < x_1x_3$}{$x_2^2 > x_1x_3$}
        \begin{jieda}
            设$x_1+y_1 = 2m$，故可设$\left \{
    \begin{array}{lll}
         x_1 = m - a, & y_1 = m + a, & a > 0\\
         x_2 = m - b, & y_2 = m + b, & b > 0\\
         x_3 = m - c, & y_3 = m + c, & c > 0
    \end{array}
    \right.$  \\
    由此，题目中的条件可以转化为 $m^2 - a^2 + m^2 - c^2 = 2(m^2-b^2)$，即$a^2+c^2 = 2b^2$，且有$m^2 - b^2 > 0$。下面推导$a+c<2b$
    \begin{enumerate}
        \item 方法1：利用平方差公式 \\
        原式移项后可以进一步可以转化为$(a+b)(a-b) + (c+b)(c-b) = 0$。\\
        $(a+b)(a-b) + (c+b)(c-b) = 0$说明$a < b < c$或者$c < b < a$，不妨设$a < b < c$（另一种情况对称），此时必有$(a+b) < (c+b)$，因此有$|a-b| > |c-b|$，即$b-a > c-b$，故$a+c<2b$
        \item 方法2：利用平均值不等式 \\
        由平均值不等式有$\sqrt{\dfrac{a^2+c^2}{2}} > \dfrac{a+c}{2}$，根据$a^2+c^2 = 2b^2$，可以转化为$2b > a + c$
    \end{enumerate}
    \end{jieda} 
    \end{enumerate}
    \item 有限制条件利用平均值不等式求最值
    \begin{enumerate}
        \item 已知$a > 0, b > 0$，且$ab = 1$，则$\frac{1}{2a} + \frac{1}{2b} + \frac{8}{a+b}$的最小值为\blkda{4}.
        \begin{jieda}
            \begin{enumerate}
                \item 消元法 \\
                用$a$表示$b$有$b = \frac{1}{a}$。原式可以化为$\frac{a}{2} + \frac{1}{2a} + \frac{8}{\frac{1}{a} + a} = \frac{1+a^2}{2a} + \frac{8a}{1+a^2} \geq 2\sqrt{\frac{1+a^2}{2a} \cdot \frac{8a}{1+a^2}} = 4$，当且仅当$a+b = a + \frac{1}{a} = 4$时取等号。 
                \item 乘常数 \\
                原式$=(\frac{1}{2a} + \frac{1}{2b}) \times ab + \frac{8}{a+b} = \frac{a+b}{2} + \frac{8}{a+b} \geq 2\sqrt{\frac{a+b}{2} \cdot \frac{8}{a+b}} = 4$，当且仅当$\frac{a+b}{2} = \frac{8}{a+b}$，即$a+b=4$时取等号。
            \end{enumerate}
            \ps 容易出现的错解：消元后分别利用均值不等式得到答案5。\textbf{应用多次平均值不等式时应注意不等式方向是否相同，取等条件是否一致。}
        \end{jieda}
        \item 
        \begin{enumerate}
            \item 已知正实数$a, b$满足$a+b = 4$，则$\frac{1}{a+1} + \frac{1}{b+3}$的最小值是\blkda{$\frac{1}{2}$}.
            \begin{jieda}
                $\frac{1}{a+1} + \frac{1}{b+3} = (\frac{1}{a+1} + \frac{1}{b+3}) \times (a+1+b+3) \times \frac{1}{8} = (2+\frac{b+3}{a+1} + \frac{a+1}{b+3}) \times \frac{1}{8} \geq \frac{1}{2}$，当且仅当$a = 3, b = 1$时取等号.
            \end{jieda}
            \item 已知正实数$a, b$满足$4a+3b = 1$，则$\frac{1}{2a+b} + \frac{1}{a+b}$的最小值是\blkda{$3+2\sqrt{2}$}.
            \begin{jieda}
                $\frac{1}{2a+b} + \frac{1}{a+b} = \frac{1}{2a+b} + \frac{2}{2a+2b} = (\frac{1}{2a+b} + \frac{2}{2a+2b}) \times (2a + b + 2a + 2b) = 1 + 2 + \frac{2a+2b}{a+b} + \frac{2(2a+b)}{2a+2b} \geq 3+2\sqrt{2}$，当且仅当$a = \frac{3\sqrt{2}}{2}-2, b = 3-2\sqrt{2}$时取等号. \\
                \ps 体会“基”的思想：将分母看做一组“基”，线性组合表示限制条件，根据系数对目标式进行变形.
            \end{jieda}
        \end{enumerate}
    \end{enumerate}
    \item 高斯取整函数
    \begin{enumerate}
        \item 记不超过$x$的最大整数为$[x]$
        \begin{enumerate}[1.]
            \item 解下列两个方程$[x] = 3, [2x] = -3$ 
            \item 设方程$[|x| + |x-1|] = 3$的解集为$A$，集合$B = \{x | 2x^2 - 11kx+15k^2 \geq 0\}$，$A \cup B = \mathbf{R}$，求实数$k$的取值范围
            \item 求方程$4x^2 - 40[x] + 51 = 0$的实数解
        \end{enumerate}
        \begin{jieda}
            \begin{enumerate}[1.]
                \item $[x] = 3$即$x \in [3, 4)$ \\
                $[2x] = -3$即$2x \in [-3, -2)$，故$x \in [-\frac{3}{2}, -1)$
                \item $[|x| + |x-1|] = 3$即$|x| + |x-1| \in [3, 4)$，故$A = [2, 2.5) \cup (-1.5, -1]$，$B = \{x | (2x-5k)(x-3k) \geq 0\}$
                \begin{dingautolist}{172}
                    \item 当$k > 0$时, $B = (-\infty, \frac{5}{2}k] \cup [3k, +\infty)$，因为$A \cup B = \mathbf{R}$，故必有$2 \leq \frac{5}{2}k, 3k \leq 2.5$，解得$k \in [\frac{4}{5}, \frac{5}{6}]$
                    \item 当$k = 0$时，B = $\mathbf{R}$，满足题意
                    \item 当$k < 0$时, $B = (-\infty, 3k] \cup [\frac{5}{2}k, +\infty)$，因为$A \cup B = \mathbf{R}$，故必有$-1 \geq \frac{5}{2}k, 3k \geq -1.5$，解得$k \in [-\frac{1}{2}, -\frac{2}{5}]$
                \end{dingautolist}
                综上，$k \in [-\frac{1}{2}, -\frac{2}{5}] \cup [\frac{4}{5}, \frac{5}{6}] \cup \{0\}$
                \item 设$x = [x] + m (m \in [0, 1))$方程可以变形为$4x^2 -40(x-m)+ 51 = 0$，进一步变化有$4x^2 -40x+ 51 = -40m$，故$4x^2 -40x+ 51 \in (-40, 0]$，解得$x \in [\frac{3}{2}, \frac{7}{2}) \cup (\frac{13}{2}, \frac{17}{2}]$故$[x] \in \{1, 2, 3, 6, 7, 8\}$，分别求解，解得合理的$x \in \{\frac{\sqrt{29}}{2}, \frac{3\sqrt{21}}{2}, \frac{\sqrt{229}}{2}, \frac{\sqrt{269}}{2}\}$ \\
                \ps \textbf{划分范围后分类讨论是解决含有高斯取整函数的方程的常规思路.}
            \end{enumerate}
        \end{jieda}
        \item 记不超过$x$的最大整数为$[x]$，若集合$P = \{x | x^2 - 2[x] - 3 = 0\}$，$Q = \{x | \frac{1}{8} \leq 2^x \leq 8\}$，则$P \cap Q = $\blkda[2.5]{$\{-1, \sqrt{7}, 3\}$}
        \begin{jieda}
            首先求得$Q = [-3, 3]$，因此$\{[x] | x \in P \cap Q \} \subseteq \{-3, -2, -1, 0, 1, 2, 3\}$，将$[x]$的值带入$P$的方程中，分别求解可得$P \cap Q = \{-1, \sqrt{7}, 3\}$
        \end{jieda}
        \item 记不超过$x$的最大整数为$[x]$，$\{x\} = x - [x]$
        \begin{enumerate*}[(a)]
            \item $[x] + [y] \leq [x+y]$
            \item $\exists x_0 \in \mathbf{R}, [\{x_0\}] \neq 0$
            \item $|x-y| < 1$是$[x] = [y]$的充分不必要条件
            \item 方程$2x\{x\}-x-1 = 0$的所有实根之和为-1.
        \end{enumerate*}
        上述命题中正确的序号是\blkda{(a), (d)}
        \begin{jieda}
            容易判断前两个分别是真和假，第三个应为必要不充分，第四个可以分离$\{x\}$，得到$2\{x\} = 1 + \frac{1}{x}$，转化为两函数的交点，根据图象和对称性可知第四个是正确的.
        \end{jieda}
    \end{enumerate}
    \item 抽象函数、复合函数、函数运算的定义域
    \begin{enumerate}
        \item 已知$f(x)=x+2, x\in[1,4]$, 则$y=f(x^2)+f^2(x)$的最大值是\blkda{22}. 
        \begin{jieda}
        先看定义域: $f(x)$的定义域为$[1,4]$，因此$f(x^2)$中$x^2\in[1,4]$, 故$f(x^2)$的定义域为$x\in[-2,-1] \cup [1,2]$. 而$f^2(x)$的定义域与$f(x)$的定义域相同.\\
        两者相加，定义域取交集，因此$y=f(x^2)+f^2(x)$的定义域为$[1, 2]$ \\
        于是$y=x^2+2+(x+2)^2=2x^2+4x+6=2(x+1)^2+4\in[12,22]$.
        \end{jieda}
        \item 已知$y = f(2x+1)$的定义域为$(1, 3]$，则$y = f(x+1)$的定义域为\blkda{$(2, 6]$}. 
        \begin{jieda}
            因为$y = f(2x+1)$的定义域为$(1, 3]$，所以$y = f(x)$的定义域是$(3,7]$，故$y = f(x+1)$的定义域为$(2, 6]$. \\
            \ps 此类题目的桥梁是\textbf{$f$括号内的取值范围是一致的}.
        \end{jieda}
        \item 下列两组函数中表示同一函数的有\blkda{\ding{172}}. \\
        \begin{enumerate*}
            \item $y = \frac{\sqrt{1-x^2}}{|x+2|}$ 和 $y = \frac{\sqrt{1-x^2}}{x+2}$
            \item $y = \sqrt{x-1} \cdot \sqrt{x-2}$和$y = \sqrt{x^2-3x+2}$
        \end{enumerate*}
        \begin{jieda}
            本质上还是在考察定义域.
        \end{jieda}
        \item 函数$y=\sqrt{\log_{\frac{1}{3}}\left(\log_{\frac{1}{3}}\left(\log_{\frac{1}{3}}x\right)\right)}$的定义域是\blkda{$\left(\frac{1}{3},\sqrt[3]{\frac{1}{3}}\right]$}.
        \begin{jieda}
        $\log_{\frac{1}{3}}\left(\log_{\frac{1}{3}}\left(\log_{\frac{1}{3}}x\right)\right)\geq 0\iff 0<\log_{\frac{1}{3}}\left(\log_{\frac{1}{3}}x\right)\leq 1
        \iff 1>\log_{\frac{1}{3}}x\geq\frac{1}{3}\iff (\frac{1}{3})^1<x\leq(\frac{1}{3})^{\frac{1}{3}}$. 
        \end{jieda}
    \end{enumerate}
    \item 根据值域寻找定义域
    \begin{enumerate}
        \item 若函数$f(x)=x^2-3x-4$的定义域为$[0,m]$, 值域为$\left[-\frac{25}{4},-4\right]$, 则实数$m$的取值范围是\blkda{$[\frac{3}{2},3]$}. 
        \begin{jieda}
        $f(x)=x^2-3x-4=(x-\frac{3}{2})^2-\frac{25}{4}$，因此$x = \frac{3}{2}$是函数值$-\frac{25}{4}$对应的唯一自变量，必然要在定义域内. 又因为$f(0)=f(3)=-4$, 是$f(x)=-4$唯二的解. 于是由二次函数的图像易得$m\in[\frac{3}{2},3]$. 
        \end{jieda}
        \item 若一系列函数的解析式相同, 值域相同, 但其定义域不同, 则称这些函数为``同族函数'', 
        那么函数解析式为$f(x)=x^2$, 值域为$\{1,4\}$的``同族函数''共有\blkda{$9$}个.
        \begin{jieda}
        $1,-1,\pm1\to 1$, $2,-2,\pm2\to 4$, 于是从值域的角度, $x$的组合共有$3\times3=9$种选择. \\
        \ps 关键是要清楚一个$y$对应几个$x$.
        \end{jieda} 
        \item 设函数$f(x)$满足$f(x) = f(\frac{1}{1+x})$对任意$x \in [0, +\infty)$都成立，其值域是$A_f$，已知对任何满足上述条件的$f(x)$都有$\{y| y = f(x), 0 \leq x \leq a \} = A_f$，则$a$的取值范围为 \blkda[2]{$[\frac{\sqrt{5}-1}{2}, +\infty)$}
        \begin{jieda}
            易知$y = x$与$y = \frac{1}{1+x}$在$x \in [0, +\infty)$只有一个交点$(\frac{\sqrt{5}-1}{2}, \frac{\sqrt{5}-1}{2})$，故在$[0,\frac{\sqrt{5}-1}{2})$和$(\frac{\sqrt{5}-1}{2}, +\infty)$内$f(x)$所对应的值域是相同的。\\
            $\forall x \in (\frac{\sqrt{5}-1}{2}, +\infty) , \exists t = \frac{1}{x+1} \in [0,\frac{\sqrt{5}-1}{2}), \quad s.t. f(x) = f(t)$ \\
            $\forall x \in [0,\frac{\sqrt{5}-1}{2}), \exists t = \frac{1}{x} - 1 \in (\frac{\sqrt{5}-1}{2}, +\infty), \quad  s.t. f(x) = f(t)$ \\
            但是$f(\frac{\sqrt{5}-1}{2})$可以和其他任意的函数值不相等。因此$a \geq \frac{\sqrt{5}-1}{2}$时满足$[0, a]$上$f(x)$的值域也是$A_f$
        \end{jieda}
    \end{enumerate}
    \item 分式型求值域
    \begin{enumerate}
        \item 求下列函数的值域: 
        \begin{enumerate*}
        \item $y=\frac{3+x}{4-x}$; 
        \item $y=\frac{2x^2+3x+8}{x^2+x+4}$; 
        \item $y=\frac{x^4+x^2+5}{(x^2+1)^2}$.
        \end{enumerate*}
        \begin{jieda}
        \begin{enumerate}
        \item $y=-1+\frac{7}{4-x}\in(-\infty,-1)\cup(-1,+\infty)$. 
        
        \item $y=2+\frac{x}{x^2+x+4}$. \\
        当$x = 0$时，$y = 2$；当$x \neq 0$时，$y = 2 + \frac{1}{x + \frac{4}{x} + 1}$.\\
        $x + \frac{4}{x} + 1 \in (-\infty, -3]\cup [5, +\infty)$，故$y \in [\frac{5}{3}, \frac{11}{5}]$
        
        \item $y=\frac{(x^2+1)^2-(x^2+1)+5}{(x^2+1)^2}=1 - \frac{1}{(x^2+1)} + \frac{5}{(x^2+1)^2}$. \\
        令$t = \frac{1}{(x^2+1)}, t \in (0, 1]$，故$y = 5t^2-t+1 \in [\frac{19}{20}, 5]$
        \end{enumerate}
        \end{jieda}
    \end{enumerate}
    \item 根式型求值域
    \begin{enumerate}
        \item 求下列函数的值域: \\
        \begin{enumerate*}
        \item $y=x+4\sqrt{1-x}$; 
        \item $y=\frac{\sqrt{x^2+1}}{x-1}$. 
        \item $y = \sqrt{1+x} - \sqrt{2+x}$.
        \item $y = \sqrt{1+x} + \sqrt{2-x}$.
        \end{enumerate*}
        \begin{jieda}
        \begin{enumerate}
        \item 令$t=\sqrt{1-x}\in[0,+\infty)$, 则$x=1-t^2$, 故$y=1-t^2+4t=-(t-2)^2+5\in(-\infty,5]$. \\
        \ps 可化为二次的. 
        
        \item 当$x>1$时, $y=\sqrt{\frac{x^2+1}{x^2-2x+1}}=\sqrt{1+\frac{2x}{x^2-2x+1}}=\sqrt{1+\frac{2}{x+\frac{1}{x}-2}}\in(1,+\infty)$. \\
        当$x<1$时, $y=-\sqrt{\frac{x^2+1}{x^2-2x+1}}=-\sqrt{1+\frac{2x}{x^2-2x+1}}$. \\
        再分为$(-\infty,0)\cup\{0\}\cup(0,1)$讨论可得$y\in(-\infty,-\frac{1}{\sqrt{2}}]$. \\
        综上, $y\in(-\infty,-\frac{1}{\sqrt{2}}]\cup(1,+\infty)$.
        \item $y = \sqrt{1+x} - \sqrt{2+x} = \frac{(\sqrt{1+x} + \sqrt{2+x})(\sqrt{1+x} - \sqrt{2+x})}{(\sqrt{1+x} + \sqrt{2+x})} = \frac{-1}{(\sqrt{1+x} + \sqrt{2+x})}$ \\
        易知函数在定义域$[-1, +\infty)$内单调递增，最小值在左端点$x = -1$取$y = -1$，当$x$趋于无穷时，函数值趋于0。因此值域为$[-1, 0)$
        \item $y = \sqrt{1+x} + \sqrt{2-x}$，故$y^2 = (1+x) + (2-x) + 2\sqrt{(1+x)(2-x)}  = 3 + 2\sqrt{(1+x)(2-x)} \in [3, 6]$，故$y \in [\sqrt{3}, \sqrt{6}]$
        \end{enumerate}
        \end{jieda}
    \end{enumerate}
    \item 复合函数的单调性
    \begin{enumerate}
        \item 函数$y = \frac{1}{x^2+2x-4}$的严格增区间是\blkda[5.5]{$(-\infty,-1-\sqrt{5})$和$(-1-\sqrt{5}, -1]$}. 
        \begin{jieda}
            \textbf{定义域先行}：$x \in \{x | x \neq -1\pm \sqrt{5}\}$ \\
            外层函数为反比例函数，有两个严格减区间，故要寻找整体的严格增区间，应寻找内层函数的严格减区间即$(-\infty,-1]$ \\
            综合两者，因此是两个单调区间：$(-\infty,-1-\sqrt{5})$和$(-1-\sqrt{5}, -1]$ \\
            \ps 单调区间之间不能用并集符号连接！
        \end{jieda}
        \item 函数$f(x) = (\frac{1}{2})^{\sqrt{x^2-3x-10}}$的递增区间是\blkda[3]{$(-\infty, -2]$}. 
        \begin{jieda}
            \textbf{定义域先行}：$x \in (-\infty, -2] \cup [5, +\infty)$ \\
            外层函数为单调减的指数函数，故要寻找整体的严格增区间，应寻找内层函数的严格减区间即$(-\infty,-2]$
        \end{jieda}
        \item 函数$f(x)=\log_2\left(x-\frac{a}{x}\right)(a>0)$在区间$[2,+\infty)$上是增函数, 则$a$的取值范围是\blkda{$(0,4)$}. 
        \begin{jieda}
        当$a>0$时, $y=x-\frac{a}{x}$在$(0,+\infty)$上单调递增. \\
        考虑到\textbf{定义域}, 于是$2-\frac{a}{2}>0$, 故$a\in(0,4)$. 
        \end{jieda}
        \item 已知函数 $f(x)=\frac{\sqrt{2021-a x}}{a-1}$ 在 $[0,1]$ 上严格减，那么实数 $a$ 的取值的范围是\blkda[3]{$(-\infty, 0) \cup(1,2021]$}.
        \begin{jieda}
        首先判断“定义域”，$2021-ax \geq 0$在$x \in [0, 1]$恒成立，设$g(x) = 2021-ax$，这是一条斜率不确定的线段，只需两个端点满足即可恒成立，即$g(0) \geq 0, g(1) \geq 0$，因此$a \leq 2021$，又因为$a \neq 1$，因此$a \in (-\infty, 1) \cup (1, 2021]$时$f(x)$在$[0, 1]$上有定义. \\
        当 $a<0$ 时, $y=\sqrt{2021-a x}$ 在 $[0,1]$ 上严格单调增, 且 $a-1<0$,所以 $f(x)$ 在 $[0,1]$ 上严格单调减, 符合题意, 当 $a=0$ 时, $f(x)=-\sqrt{2021}$ 并非严格增或严格减, 不符合题意, 当 $0<a<1$ 时, $y=\sqrt{2021-a x}$ 在 $[0,1]$ 上严格单调减, 且 $a-1<0$, 不符合题意, 当 $2021\geq a>1$ 时, $y=\sqrt{2021-a x}$ 在 $[0,1]$ 上严格单调减, $a-1>0$, 符合题意.\\
        综上实数 $a$ 的取值的范围是 $(-\infty, 0) \cup(1,2021]$. \\
        \ps \textbf{$y = a \cdot f(x)$，即一个函数乘常系数，类型的函数的单调性在$a > 0$时与$f(x)$保持一致，在$a < 0$时与$f(x)$完全相反.}
        \end{jieda}
        \item 若函数$y=\frac{a^x}{2a^2-5a+2}$在$\bR$上为严格减函数, 则$a$的取值范围是\blkda[3]{$(0,\frac{1}{2})\cup(1,2)$}. 
        \begin{jieda}
        $\begin{cases}
        a\in(0,1)\\
        2a^2-5a+2>0
        \end{cases}$
        或$\begin{cases}
        a>1\\
        2a^2-5a+2<0
        \end{cases}$, 
        解得$a\in(0,\frac{1}{2})\cup(1,2)$. 
        \end{jieda}
    \end{enumerate}
    \item 割线斜率与单调性
    \begin{enumerate}
        \item 定义在 $(0,+\infty)$ 上的函数 $f(x)$ 满足 $\frac{x_{2} f\left(x_{1}\right)-x_{1} f\left(x_{2}\right)}{x_{1}-x_{2}}<0$, 且 $f(3)=9$, 则使 $f(x)>3 x$ 成立的 $x$ 的取值范围是\blkda{$(0,3)$} .
        \begin{jieda}
        由 $\frac{x_{2} f\left(x_{1}\right)-x_{1} f\left(x_{2}\right)}{x_{1}-x_{2}}<0$, 且 $x_{1}, x_{2} \in(0,+\infty)$, 则两边同时除以 $x_{1} x_{2}$ 可得 $\frac{\frac{f\left(x_{1}\right)}{x_{1}}-\frac{f\left(x_{2}\right)}{x_{2}}}{x_{1}-x_{2}}<0$,
        令 $g(x)=\frac{f(x)}{x}$, 则原不等式为 $\frac{g\left(x_{1}\right)-g\left(x_{2}\right)}{x_{1}-x_{2}}<0$, 因此函数 $g(x)$ 在 $(0,+\infty)$ 上严格减,由 $f(x)>3 x$, 得 $\frac{f(x)}{x}>3$, 又 $g(3)=\frac{f(3)}{3}=3$, 于是 $g(x)>g(3)$, 解得 $0<x<3$,所以使 $f(x)>3 x$ 成立的 $x$ 的取值范围是 $(0,3)$. \\
        \ps $\frac{f(x_1) - f(x_2)}{x_1 - x_2}$的几何意义是割线的斜率，遇到接近$\frac{f(x_1) - f(x_2)}{x_1 - x_2}$的情况，可以凑出割线斜率的形式，若割线斜率恒大于零则函数在对应区间上严格增.
        \end{jieda}
        \item 函数$f(x)$为定义在$(-\infty, 0) \cup (0, +\infty)$上的奇函数，$f(2) = 1$，对于任意$x_1, x_2 \in (0, +\infty), x_1 \neq x_2$，都有$\frac{x_1f(x_1) - x_2f(x_2)}{x_1 - x_2} > 0$成立，则$f(x) \leq \frac{2}{x}$的解集为\blkda[3]{$(-\infty， -2] \cup (0, 2]$}.
        \item 设奇函数$f(x)$在$(0, +\infty)$上为增函数，且$f(1) = 0$，则不等式$\frac{f(x) - f(-x)}{x} < 0$的解集为\blkda[3]{$(-1, 0) \cup (0, 1)$}.
    \end{enumerate}
    \item 分段函数严格单调的充要条件
    \begin{enumerate}
        \item 已知函数 $f(x)=\left\{\begin{array}{l}x^{2}+2 a x, x \geq-1 \\ (2-a) x+2 a-5, x<-1\end{array}\right.$, 若 $f(x)$ 在 $\mathbf{R}$ 上是严格增函数, 则实数 $a$ 的取值范围是\blkda{$1 \leq a \leq \frac{8}{5}$} .
        \begin{jieda} 由函数 $f(x)=\left\{\begin{array}{l}x^{2}+2 a x, x \geq-1 \\ (2-a) x+2 a-5, x<-1\end{array}\right.$ 在 $\mathbf{R}$ 上是严格增函数, \\
        得 $\left\{\begin{array}{l}-a \leq-1 \\ 2-a>0 \\ 1-2 a \geq-(2-a)+2 a-5\end{array}\right.$,解得 $1 \leq a \leq \frac{8}{5}$,
        所以实数 $a$ 的取值范围是 $1 \leq a \leq \frac{8}{5}$.
        故答案为: $1 \leq a \leq \frac{8}{5}$ \\
        \ps 分段函数严格单调增等价于各段都严格增，且在端点处相连或者满足“更上一层楼型”
        \end{jieda}
        \item 已知函数$f(x) = \left \{ \begin{array}{ll}
             (2-5a)^x & x < 1  \\
             \frac{a}{x} + 5 & x \geq 1
        \end{array}\right.$ 在$\mathbf{R}$上严格递增，则实数$a$的取值范围是\blkda{$[-\frac{1}{2}, 0)$} .
        \begin{jieda}
            函数$f(x)$在$\mathbf{R}$上严格递增的充要条件为$\left \{ \begin{array}{l}
             2-5a > 1  \\
             a < 0 \\
             a + 5 \geq 2 - 5a
        \end{array}\right.$，解得$a \in [-\frac{1}{2}, 0)$
        \end{jieda}
    \end{enumerate}
    \item 含参的分段函数的最值
    \begin{enumerate}
    \item 已知函数$f(x) = \left\{  
         \begin{array}{ll}  
            x^2-2ax+9, & x \leq 1, \\  
            x+\frac{4}{x}+a, & x > 1.
         \end{array}
    \right.$ 若$f(x)$的最小值为$f(1)$，则实数$a$的取值范围是\blkda{$[2, +\infty)$}
    \begin{jieda}
        首先看到函数在$x \leq 1$时是开口向上的二次函数，函数在$x > 1$时是对勾函数。由于$f(1)$是最小值，所以$f(x)$在$(-\infty, 1]$应该单调递减，故其对称轴应该位于$x=1$及其右侧，即$a \geq 1$。这样我们有$f(1) = 10-2a < f(x), x < 1$。又因为$f(x)$在$(1, +\infty)$为对勾函数，在$(1, 2]$单调递减，在$(2, +\infty)$单调递增，$f(x)$在$(1, +\infty)$的最小值在$x = 2$取得，$f(2) = 4+a$, 故$4+a \geq 10-2a$, 故$a \geq 2$ \\
        \ps 分段函数的最小值要比每一段内所有的函数值都要小.
    \end{jieda}
    \item 设$a$为实数, 函数$f(x)=2x^2+(x-a)\abs{x-a}$. \\
    \begin{enumerate*}[1.]
    \item 若$f(0)\geq 1$, 求$a$的取值范围; 
    \item 求$f(x)$的最小值. 
    \end{enumerate*}
    \begin{jieda}
    \begin{enumerate}[1.]
    \item $f(0)=-a\abs{a}$，因为$g(x)=-x\abs{x}$在$\bR$上严格单调递, 且$g(-1)=1$. \\
    于是$f(0)=-a\abs{a}\geq 1=g(-1)\iff a\leq -1$.  \\
    \ps 想象一下$g(x)=-x\abs{x}$的图象.
    \item $f(x)=\begin{cases}
    x^2+2ax-a^2, & x\leq a\\
    3x^2-2ax+a^2, & x\geq a
    \end{cases}$. \\
    函数图象由两个二次函数图象的部分拼接而成，且拼接处连续，分析两个开口向上的二次函数的对称轴可知，一个为$x = -a$，一个为$\frac{a}{3}$ \\
    讨论拼接点$x = a$与两个对称轴不同的位置关系下函数的单调性，即可得到函数的最小值，在不同的拼接点条件下寻找全局的最小值即为所求. 
    \begin{dingautolist}{172}
        \item $a = 0$时，三者重合，此时最小值为$0$ 
        \item $a > 0$时，$-a < \frac{a}{3} < a$，故函数$f(x)$在第一段上先减后增，在第二段上单调增，因此最小值在第一段的顶点处取得（$x = -a$），最小值为$-2a^2$
        \item $a < 0$时，$a < \frac{a}{3} < -a$，故函数$f(x)$在第一段上单调递减，在第二段上先减后增，因此最小值在第二段顶点取得（$x = \frac{a}{3}$），最小值为$\frac{2a^2}{3}$
    \end{dingautolist}
    综上，$a < 0$时为$\frac{2a^2}{3}$，$a \geq 0$时，$f(x)$的最小值为$-2a^2$. \\
    \ps 最值需要通过单调性来说明，分段函数分段求，分析各段上函数的单调性.
    另外, 本题的另一个麻烦之处在于函数中的分段点与参数$a$有关, 各段中的``对称轴''(单调区间的端点)亦与参数$a$有关. 
    \end{enumerate}
    \end{jieda}
    \item 已知函数$f(x) = x^2 + (x-1)|x-a|$，是否存在实数$a$，使得不等式$f(x) \geq 2x-3$对一切实数$x$恒成立，若存在，求出$a$的范围，若不存在，说明理由.
    \begin{jieda}
        不等式$f(x) \geq 2x-3$对一切实数$x$恒成立即$f(x) - 2x + 3 \geq 0$恒成立. 设$g(x) = f(x) - 2x + 3$，$g(x) = \left \{ \begin{array}{ll}
            (a-1)x+(3-a) & x < a \\
            2x^2 - (a+3)x+(a+3) & x \geq a
        \end{array} \right.$，故应有$g(x) \geq 0$恒成立.
        \begin{dingautolist}{172}
            \item $a > 1$ ，此时$g(x)$在$(-\infty, a)$为单调递增的一次函数，不满足要求.
            \item $a = 1$，此时$g(x) = 2(x < a)$，当$x \geq a$时$g(x)$单调递增，故满足要求.
            \item $a < 1$，此时$g(x)$在$(-\infty, a)$为单调递减的一次函数，在$[a, \frac{a+3}{4})$单调减，在$[\frac{a+3}{4}, +\infty)$单调增，故$g(x) \geq g(\frac{a+3}{4}) = (a+3) - \frac{(a+3)^2}{8}$，当$(a+3) - \frac{(a+3)^2}{8} \geq 0$时满足题意，解得$a \in [-3, 1)$
        \end{dingautolist}
        综上，$a \in [-3, 1]$ \\
        \ps 函数解析式内有部分绝对值的情况通常要写成分段函数来分析
    \end{jieda}
    \end{enumerate}
    \item 奇偶性概念
    \begin{enumerate}
        \item 函数$f(x)=x\abs{x+a}+b$是奇函数的充要条件是\dotfill(\kh{B})
        \xx{$ab=0$}{$a^2+b^2=0$}{$a=b$}{$a+b=0$}
        \begin{jieda}
        奇函数若在$x=0$处有定义, 则此处的函数值为零, 即$f(0)=0$, 于是$b=0$. \\
        故$f(x)=x\abs{x+a}$. 于是由$f(x)$为奇函数，不难得到$a=0$, 比如$f(-a)+f(a)=0$. 
        \end{jieda}
        \item 若函数$f(x)=\frac{a-\me^x}{1+a\cdot\me^x}$($a$为常数)在定义域上为奇函数, 则$a$的值为\blkda{$\pm 1$}. 
        \begin{jieda}
        首选解法：$f(0)=0$或无意义，即$\frac{a-1}{1+a} = 0$或$\frac{a-1}{1+a}$无意义，故$a = \pm 1$\\
        也可以采用就极限思想：让函数在$\pm \infty$时的极限互为相反数求解：$x \rightarrow +\infty, f(x) \rightarrow -\frac{1}{a}, x \rightarrow -\infty, f(x) \rightarrow a$，由$a = \frac{1}{a}$得$a = \pm 1$. \\
        把握定义求解有：$f(x) = \frac{a-\me^x}{1+a\cdot\me^x}, f(-x) = \frac{a-\me^{-x}}{1+a\cdot\me^{-x}} = \frac{a\me^{x}-1}{a+\me^{x}}$ \\
        由$f(x) = -f(-x)$知$\frac{a-\me^x}{1+a\cdot\me^x} = \frac{1-a\me^{x}}{a+\me^{x}}$，交叉相乘得到$a^2 - e^{2x} = 1-a^2e^{2x}$，移项整理得$(a^2-1) = (1-a^2)\cdot(e^{2x})$，因为此式恒成立，所以必有$a = \pm 1$
        \end{jieda}
        \item 已知$\mathbf{R}$上的奇函数与偶函数$g(x)$满足$f(x) - g(x) = 2^x$，若对于$\forall x \in (0, 2]$，$mf(x) - g(2x) \geq 0$恒成立，则$m$的取值范围是\blkda[2]{$[-2\sqrt{2}, +\infty)$}
        \begin{jieda}
            $f(x) = \frac{2^x - 2^{-x}}{2}, g(x) = -\frac{2^x+2^{-x}}{2}$ \\
            故$m(\frac{2^x - 2^{-x}}{2}) + \frac{4^x+4^{-x}}{2} \geq 0$在$(0, 2]$上恒成立. \\
            换元$t = 2^x - 2^{-x}, t \in (0, \frac{15}{4}]$，则$4^x+4^{-x} = t^2+2$，故有$mt + t^2 + 2 \geq 0$在$(0, \frac{15}{4}]$上恒成立. \\
            参变分离有$m \geq -(t + \frac{2}{t})$在$(0, \frac{15}{4}]$上恒成立.由对勾函数性质可知$-(t + \frac{2}{t})$在$(0, \frac{15}{4}]$上最大值为$-2\sqrt{2}$，因此有$m \in [-2\sqrt{2}, +\infty)$ \\
            \ps \textbf{一个定义域关于原点对称的函数$f(x)$一定能够拆成一个奇函数和一个偶函数之和：$f(x) = \frac{f(x) - f(-x)}{2} + \frac{f(x)+f(-x)}{2}$}
        \end{jieda}
        \item 已知函数$f(x) = \left \{ \begin{array}{ll}
            x^2 + x & x \leq 0 \\
            ax^2 + bx & x > 0
        \end{array}\right.$为奇函数，则$a+b = $\blkda{0}
        \begin{jieda}
            代入特值有$f(1) = a+b$，$f(-1) = 0$，因为$f(x)$是奇函数，所以有$a+b = 0$ \\
            \ps 根据奇偶性求参数可以考虑代入特值（必要条件），但最后需要检验是否充要（本题只问$a+b$并不涉及）.
        \end{jieda}
        \item 函数$f(x) = \frac{x^2+(3a+1)x+c}{x+a} (a, c \in \mathbf{R})$
        \begin{enumerate}
            \item 当$a = 0$时，是否存在实数$c$使得$f(x)$是奇函数
            \item 函数$f(x)$的图像过点$(1, 3)$，且$f(x)$的图像和$x$轴负半轴有两个交点，求实数$a$的取值范围
        \end{enumerate}
        \begin{jieda}
            \begin{enumerate}
            \item 当$a = 0$，$f(x) = \frac{x^2+x+c}{x} = x + 1 + \frac{c}{x} (x \neq 0)$， $f(1) + f(-1) = 2 \neq 0$，故不存在$c$使得$f(x)$为奇函数.
            \item $f(1) = \frac{2+3a+c}{1+a} = 3$, 故$c = 1, f(x) = \frac{x^2+(3a+1)x+1}{x+a}$. 函数的零点是函数的分子部分的零点子集，故$g(x) = x^2+(3a+1)x+1 (x \neq -a)$与$x$轴负半轴有两个交点$x_1, x_2$。
            \begin{equation*}
            \left\{ 
                \begin{array}{l}  
                    \Delta = (3a+1)^2 - 4 > 0, \\ 
                    x_1 + x_2 = -(3a + 1) < 0, \\
                     a^2 - (3a+1)a + 1 \neq 0.
                 \end{array}
            \right \quad
            \end{equation*}
            解得$a > \frac{1}{3} \mbox{且} a \neq \frac{1}{2}$
        \end{enumerate}
        \ps 否定一个函数是奇函数可以通过举反例.
        \end{jieda}
    \end{enumerate}
    \item 反比例型函数的对称中心
    \begin{enumerate}
        \item 函数$y=\frac{3-x}{1+2x}$的严格递减区间是\blkda[4]{$(-\infty,-\frac{1}{2})$和$(-\frac{1}{2},+\infty)$}, 对称中心是\blkda{$(-\frac{1}{2},-\frac{1}{2})$}, \\
        值域是\blkda[2]{$\{t\mid t\neq-\frac{1}{2}\}$}. 
        \begin{jieda}
            反比例函数对称中心的横坐标是分母为零的值，纵坐标是分离出的常数.
        \end{jieda}
        \item 函数$f(x) = \frac{x}{ax+b}$（$a, b$是非零常数），满足$f(2) = 1$，且方程$f(x) = x$仅有一个解.
        \begin{enumerate}
            \item 求$a, b$的值
            \item 是否存在常数$m$，使得对于定义域内任意的$x$，$f(x) + f(m-x) = 4$恒成立
            \item 在平面直角坐标系中，求定点$A(-3, 1)$到此函数上任意一点$P$的距离$|AP|$的最小值.
        \end{enumerate}
        \begin{jieda}
            \begin{enumerate*}
                \item 由$f(2) = 1$，$\frac{2}{2a+b} = 1$，故$2a+b = 2$；$f(x) = x$即$\frac{x}{ax+b} = x$，即$(\frac{1}{ax+b}-1)x = 0$，容易判断其有两个解$x_1  = 0, x_2 = \frac{1-b}{a}$.要符合题意只可能是$b = 1$，两根重合.进一步推得$a = \frac{1}{2}$ \\
                \item $f(x) = \frac{2x}{x+2}$，为反比例型函数，其对称中心为$(-2, 2)$，因此$f(x)$满足$f(x) + f(-4-x) = 4$. 由此可知$m = -4$ \\
                \item 设$P(x_0, \frac{2x_0}{x_0+2})$，故$|AP| = \sqrt{(x_0+3)^2+(\frac{2x_0}{x_0+2}-1)^2}   = \sqrt{(x_0+3)^2+(1-\frac{4}{x_0+2})^2}$. 设$t = x_0+2 (t \neq 0)$，则$|AP| = \sqrt{(t+1)^2+(1-\frac{4}{t})^2} = \sqrt{(t-\frac{4}{t})^2+2(t-\frac{4}{t})+10}$，设$k = t - \frac{4}{t}, k \in \mathbf{R}$，$|AP| = \sqrt{k^2+2k+10} = \sqrt{(k+1)^2+9} \geq 3$（当$x = \frac{-5\pm\sqrt{17}}{2}$时取得最小值）
            \end{enumerate*}
        \end{jieda}
    \end{enumerate}
    \item 利用函数性质解不等式
    \begin{enumerate}
        \item 定义在$\bR$上的偶函数$f(x)$满足: 对任意的$x_1,x_2\in[0,+\infty)$, $x_1\neq x_2$, 有$\frac{f(x_2)-f(x_1)}{x_2-x_1}<0$, 则\dotfill(\kh{A})
        \xx{$f(3)<f(-2)<f(1)$}{$f(1)<f(-2)<f(3)$}{$f(-2)<f(1)<f(3)$}{$f(3)<f(1)<f(-2)$}
        \begin{jieda}
        等价定义: 
        \begin{itemize}
        \item 严格单调增: $\frac{f(x_2)-f(x_1)}{x_2-x_1}>0(x_1\neq x_2)$; 
        \item 严格单调减: $\frac{f(x_2)-f(x_1)}{x_2-x_1}<0(x_1\neq x_2)$. 
        \end{itemize}
        \end{jieda}
                \item 已知函数$f(x) = -x^3-x$，若$f(t^2+2t+k)+f(-2t^2+2t-5)>0$对任意的$t \in [0, 5]$都成立，则$k$的取值范围是\blkda{$(-\infty, 1)$}.
    \begin{jieda}
        $f(x)$为奇函数，在$\mathbf{R}$上单调递减，根据对称性可以将函数值上的不等关系转化为变量的不等关系，即$(t^2+2t+k)+(-2t^2+2t-5)<0$在$t \in [0, 5]$恒成立.\\
        参变分离有$-t^2+4t-5 < -k$在$t \in [0, 5]$恒成立，因此有$k < 1$ \\
        \ps $f(x)$关于$(a, b)$对称，则$f(x_1) + f(x_2) = 2b \Leftrightarrow x_1 + x_2 = 2a$ \\
        $f(x)$关于$(a, b)$对称，且在定义域内单调增，则$f(x_1) + f(x_2) < 2b \Leftrightarrow x_1 + x_2 < 2a$ \\
        $f(x)$关于$(a, b)$对称，且在定义域内单调增，则$f(x_1) + f(x_2) > 2b \Leftrightarrow x_1 + x_2 > 2a$ \\
        $f(x)$关于$(a, b)$对称，且在定义域内单调减，则$f(x_1) + f(x_2) > 2b \Leftrightarrow x_1 + x_2 < 2a$ \\
        $f(x)$关于$(a, b)$对称，且在定义域内单调减，则$f(x_1) + f(x_2) < 2b \Leftrightarrow x_1 + x_2 > 2a$ \\
        以上结论可以通过代数推导，更推荐结合图像理解.
    \end{jieda}
    \item 已知函数$f(x) = 2^x - (\frac{1}{2})^x+3$，且$f(m^2)+f(-5m+4)<6$，则实数$m$的取值范围是\blkda{$(1, 4)$}.
    \begin{jieda}
        直接应用上面结论即可，也可以设$g(x) = f(x) - 3$，将$f(m^2)+f(-5m+4)<6$转化为$g(m^2)+g(-5m+4)<0$后按照更熟悉的奇函数形式处理.
    \end{jieda}
    \item 已知函数$y = f(3x+1)$为偶函数，且在$[0, +\infty)$上为严格增函数，若$f(x) < f(2x+1)$，则$x$的取值范围是\blkda[5]{$(-\infty, -1) \cup (\frac{1}{3}, +\infty)$}.
    \begin{jieda}
        函数$y = f(3x+1)$为偶函数，且在$[0, +\infty)$上为严格增函数即$f(x)$关于$x = 1$对称，且“开口向上”，因此离对称轴$x = 1$近的变量的函数值小，由此$f(x) < f(2x+1)$可以转化为$|x-1| < |2x|$，两侧平方有$(x-1)^2 < 4x^2$，解得$x \in (-\infty, -1) \cup (\frac{1}{3}, +\infty)$
    \end{jieda}
    \item 已知函数$f(x)$在定义域$[2-a, 3]$上是偶函数，在$[0, 3]$上严格单调递减，并且$f(-m^2-\frac{a}{5}) > f(-m^2+2m-2)$，则$m$的取值范围是\blkda[3]{$[1-\sqrt{2}, \frac{1}{2}]$}.
    \begin{jieda}
        首先根据定义域的对称性可以求得$a = 5$，进一步根据对称性可以判断出函数“开口向下”，离$y$轴越近函数值越大，再结合定义域的要求，则$f(-m^2-\frac{a}{5}) > f(-m^2+2m-2) \Leftrightarrow \left \{ \begin{array}{l}
             |-m^2-1| < |-m^2+2m-2|  \\
             -3 \leq -m^2-1 \leq 3  \\
             -3 \leq -m^2+2m-2 \leq 3 
        \end{array}\right.$第一个配方后可以直接去绝对值，最终得到$m \in [1-\sqrt{2}, \frac{1}{2}]$ \\
        \ps \textbf{奇函数和偶函数的定义域都是关于原点对称的，若要证明奇偶性首先要判断定义域！}
    \end{jieda}
    \end{enumerate}
    \item 利用函数的对称性解决问题
    \begin{enumerate}
        \item 已知函数$f(x)=x^5+ax^3+bx-8$, 且$f(-2)=10$, 则$f(2)$=\blkda{$-26$}
        \begin{jieda}
        $f(x)+f(-x)=-16$，函数关于$(0, -8)$中心对称。\\
        \ps 设$g(x) = f(x)+8$，则$g(x)$为奇函数. (多项式函数是奇函数等价于偶次方的系数为零, 是偶函数等价于奇次方的系数为零.)
        \end{jieda}
        \item 已知$f(x) = \frac{(2^x+1)^2}{2^x \cdot x}+1$在$[-2018, -1]\cup[1, 2018]$上最大值为$M$,最小值为$m$,则$M+m = $\blkda{$2$}.
        \begin{jieda}
        设$g(x) = \frac{(2^x+1)^2}{2^x \cdot x}$，则$g(x) = \frac{2^x+\frac{1}{2^x}+2}{x} = \frac{2^x+2^{-x}+2}{x}$，故$g(-x) = -g(x)$，即$g(x)$为奇函数，故$M+m = 2$ \\
        \ps \textbf{涉及最大值与最小值和的题目一般可以优先考虑对称性}.
        \end{jieda}
        \item 已知函数$f(x) = \frac{3^{x+1}+1}{3^x+1}$在$[-2021, 2021]$上的最大值和最小值分别为$M, m$，则$M+m = $\blkda{$4$}.
        \begin{jieda}
            $g(x) = f(x) - 2 = \frac{3^x - 1}{3^x+1}$为奇函数，故$f(x)$关于$(0, 2)$中心对称，$M+m = 4$ \\
            \ps 本题如果没有发现$f(x) - 2$是奇函数，利用单调性也是可以解决的，计算量相对大. \\
            \ps 其实$y=\frac{a^x-1}{a^x+1}$, $y=\frac{a^x+1}{a^x-1}$都是奇函数.
        \end{jieda}
        \item 设函数$f(x)$的定义域为$\mathbf{R}$，$f(x+1)$是奇函数，$f(x+2)$是偶函数，当$x \in [1, 2]$时，$f(x) = kx+b$，若$f(0) + f(3) = 8$，则$f(\frac{2021}{2}) = $\blkda{-4}
        \begin{jieda}
        因为$f(x)$关于$(1, 0)$中心对称，当$x \in [1, 2]$时，$f(x) = kx+b$为连续函数，所以$f(1) = 0$ 且在$[0, 2]$上$f(x) = kx+b$. 又因为$f(x)$关于$x = 2$对称，所以$f(3) = f(1)$, 故$f(0) = 8, f(\frac{1}{2}) = 4, f(\frac{3}{2}) = -4$，即$f(\frac{8k+1}{2}) = 4, f(\frac{8k+3}{2}) = -4, f(\frac{8k+5}{2}) = -4, f(\frac{8k+7}{2}) = 4(k \in \mathbf{Z})$，因此$f(\frac{2021}{2}) = -4$ \\
        \ps 如果对称中心的横坐标在定义域内，则对称中心在函数图象上，如果对称中心横坐标不在定义域内，则对称中心不在函数图象上；本题可以得到周期性的结论，但是对于解决问题并没有直接帮助。
        \end{jieda}
        \item 已知函数$g(x) = x^3 - 9x^2 + 29x - 30$，$g(m) = -12, g(n) = 18$，则$m+n = $\blkda{}
        \begin{jieda}
            对于三次函数$y = ax^3+bx^2+cx+d$，当$b^2 - 3ac \leq 0$时，其为单调函数，因此$g(x)$是单调函数. \\
            对于三次函数$y = ax^3+bx^2+cx+d$，其对称中心为$(-\frac{b}{3a}, f(-\frac{b}{3a}))$，故$g(x)$对称中心为$(3, 3)$，因为$\frac{g(m) + g(n)}{2} = 3$，所以$m+n = 6$
        \end{jieda}
        \item 已知函数$f(x) = \left \{ \begin{array}{ll}
            2 - |x| & x \leq 2 \\
            (x-2)^2 & x > 2
        \end{array}\right.$，函数$g(x) = b - f(2-x)$，如果$y = f(x) - g(x)$恰有两个零点, 则实数$b$的取值范围是\blkda[3]{$\{\frac{7}{4}\} \cup (2, +\infty)$}
        \begin{jieda}
            设$h(x) = f(x) - g(x) = f(x) + f(2-x) - b$，$h(x)$的零点即$y = f(x) + f(2-x)$与$y = b$的交点的横坐标. 因为$f(x)$与$f(2-x)$关于$x = 1$对称，所以$y = f(x) + f(2-x)$关于$x = 1$对称. 当$x \in [0, 2]$时，$f(x) + f(2-x) = 2$；当$x > 2$或$x < 0$时，$f(x) + f(2-x)$为二次函数，最小值为$\dfrac{7}{4}$ \\
            根据对称性绘制图像可知当$b \in \{\frac{7}{4}\} \cup (2, +\infty)$时满足要求. \\
            \ps \textbf{两个互对称的函数相加的结果会具有自对称性.}
        \end{jieda}
    \end{enumerate}
    \item 迭代推出周期性
    \begin{enumerate}
        \item 已知定义在$\mathbf{R}$上的偶函数$f(x)$满足$f(x) = -f(\pi - x)$，当$x \in [0, \frac{\pi}{2})$时$f(x) = \sqrt{x}$，则函数$y = f(x)$在$[-\frac{3\pi}{2}, \frac{5\pi}{2}]$上的零点个数为\blkda{9}
        \begin{jieda}
            $f(x) = -f(\pi - x)$即$f(x)$关于$(\frac{\pi}{2}, 0)$中心对称，又因为$f(x)$是偶函数，所以$2\pi$是$f(x)$的一个周期. 因为定义域为$\mathbf{R}$，所以$f(\frac{k\pi}{2}) = 0, k \in \mathbf{Z}$；结合图像可知有9个零点 \\
        \ps 两个横坐标不同位置的对称性会产生周期性
        \begin{enumerate}
            \item 两个对称轴的横坐标之差的二倍是函数的周期：$f(x) = f(2a-x), f(x) = f(2b - x) \Leftrightarrow 2|a-b|$是$f(x)$的一个周期。
            \item 两个纵坐标相同的对称中心的横坐标之差的二倍是函数的周期：$f(x) + f(2a-x) = 2m, f(x) + f(2b - x) = 2m \Leftrightarrow 2|a-b|$是$f(x)$的一个周期。
            \item 一个对称轴和一个对称中心的横坐标之差的四倍是函数的周期：$f(x) = f(2a-x), f(x) + f(2b - x) = 2m \Leftrightarrow 4|a-b|$是$f(x)$的一个周期。
        \end{enumerate}
        \end{jieda}
        \item 已知函数$f(x)$的定义域为$\mathbf{R}$，$f(x+y)+f(x-y) = f(x) \cdot f(y)$，$f(1) = 0$，则$f(2022) = $\blkda{$-2 \mbox{或} 0$}
        \begin{jieda}
            令$y = 1$，则对于任意的$x \in \mathbf{R}$有$f(x+1) + f(x-1) = 0$ \\ 故迭代有$f(x+3) + f(x+1) = 0$，因此有$f(x + 3) = f(x-1)$，即$4$是$f(x)$的一个周期. 故$f(2022) = f(2)$\\
            代入$x = 0, y = 0$，有$f(0) + f(0) = f(0) \cdot f(0)$，故$f(0) = 0$或$f(0) = 2$ \\
            代入$x = 1, y = 1$，由$f(2) + f(0) = f(1) \cdot f(1)$，故$f(2) = -f(0)$ \\
            因此$f(2022) = -2 \mbox{或} 0$
        \end{jieda}
        \item 已知函数$f(x)\neq -1$, 且对定义域内任意的$x$总有关系$(f(x+\pi)+1)(f(x)+1)=2$, 那么下列结论中正确的是\dotfill(\kh{B})
        \xx{$f(x)$是周期为$\pi$的周期函数}
        {$f(x)$是周期为$2\pi$的周期函数}
        {$f(x)$是周期为$\frac{\pi}{2}$的周期函数}
        {$f(x)$不是周期函数}
        \begin{jieda}
        由条件可得$(f(x+2\pi)+1)(f(x+\pi)+1)=2$, 于是$f(x+2\pi)=f(x)$. \\
        \ps $f(x) + f(x+t) = c\Leftrightarrow 2t$是$f(x)$的一个周期。\\
        进阶版：$f(x+a) + f(x+b) = c\Leftrightarrow 2|a-b|$是$f(x)$的一个周期。 \\
        $f(x) \cdot f(x+t) = c (c \neq 0)\Leftrightarrow 2t$是$f(x)$的一个周期。\\
        进阶版：$f(x+a) \cdot f(x+b) = c (c \neq 0)\Leftrightarrow 2|a-b|$是$f(x)$的一个周期。
        \end{jieda}
        \item 已知$f(x)$是$\bR$上的函数, 且$f(x+2)(1-f(x))=1+f(x)$. 
        \begin{enumerate*}[1.]
        \item 求证: $f(x)$是周期函数;
        \item 若$f(4)=-\sqrt{3}$, 求$f(2020)$. 
        \end{enumerate*} 
        \begin{jieda}
        \begin{enumerate}[1.]
        \item 若$f(x)$为常数函数, 则设为$f(x)=t$, 于是$t(1-t)=1+t$, 无解, 故$f(x)$不是常函数. 
        \begin{dingautolist}{172}
            \item 若$f(x)=1$有解, 设为$x_0$, 则$f(x_0+2)(1-f(x_0))=1+f(x_0)$, 故$f(x_0)=-1$. 矛盾. \\
            若$f(x)=0$有解, 设为$x_0$, 则$f(x_0+2)(1-f(x_0))=1+f(x_0)$, 故$f(x_0+2)=1$. 显然不可能(见上一行).
            \item 由(i)知$f(x)=1$无解, 且$f(x)=0$无解, \\
            则$f(x+4)=\frac{1+f(x+2)}{1-f(x+2)}=\frac{1+\frac{1+f(x)}{1-f(x)}}{1-\frac{1+f(x)}{1-f(x)}}=-\frac{1}{f(x)}$, 故$f(x+8) = -\frac{1}{f(x+4)} = f(x)$.
        \end{dingautolist}
        综上, $f(x)$是周期为8的周期函数.
        \item $f(2020)=f(4)=-\sqrt{3}$. 
        \end{enumerate}
        \end{jieda}
    \end{enumerate}
    \item 其他周期性题目
    \begin{enumerate}
        \item $f(x)$是定义在$\mathbf{R}$上的以3为周期的奇函数，且$f(2) = 0$，则方程$f(x) = 0$在区间$(0, 6)$内的解的个数的最小值是\blkda{$7$}. 
        \begin{jieda}
            $f(x)$为$\mathbf{R}$上奇函数，故$f(0) = 0$，又因为$f(2) = 0$，由周期性$f(-1) = 0$，由奇函数性质进一步可以推得$f(-2) = 0$，故由周期性可知$f(n) = 0 (n \in \mathbf{Z})$ \\
            由奇函数和周期性可知$f(1.5) = f(-1,5) = 0$，进一步可知$f(4.5) = 0$ \\
            \ps 对于周期为$T$的奇函数，必有$f(\frac{T}{2}) = f(-\frac{T}{2}) = 0$
        \end{jieda}
        \item 已知$f(x)$是定义域为$(-\infty, +\infty)$的奇函数，满足$f(1-x) = f(1+x)$。若$f(1) = 2$，则$f(1)+f(2)+f(3)+...+f(50) = $\dotfill(\kh{C})
        \xx{-50}
        {0}
        {2}
        {50}
        \begin{jieda}
            因为$f(x)$为奇函数，且$f(1-x) = f(1+x)$，故$f(x)$具有周期性，最小正周期为4，即$f(x) = f(x+4)$. 由题意可知$f(0) = 0, f(1) = 2$, 可以推出$f(2) = 0, f(3) = f(-1) = -2$，由此可知$f(4k)+f(4k+1)+f(4k+2)+f(4k+3)=0, k \in \mathbf{Z}$。故原式$= f(1) + f(2) = 2$ \\
            \ps 具有周期性的求和问题，先求一个周期内的和.
        \end{jieda}
    \end{enumerate}
    \item 根据函数性质求解析式
    \begin{enumerate}
        \item 已知函数$f(x)$是周期为$4$的函数，若$x \in [-2, 2]$时，$f(x) = x^2 + x^4$，则$x \in [2, 6]$时求$f(x)$解析式.
        \begin{jieda}
            任取$x \in [2, 6]$，因为$f(x)$周期为4，所以$f(x) = f(x - 4)$，又因为$x - 4 \in [-2, 2]$，故$f(x-4) = (x-4)^2 + (x-4)^4$，因此$f(x) = (x-4)^2 + (x-4)^4 (x \in [2, 6])$ \\
            \ps 求解析式的关键点：
            \begin{dingautolist}{172}
                \item 求哪个区间就在哪个区间任取
                \item 根据函数的性质写出一个关于$f(x)$的等式，此时另一个$f$括号内的取值范围应该是给定解析式的区间
                \item 代入已知解析式，整理即可
                \item 最后将完整的解析式（包括定义域）呈现出来.
            \end{dingautolist}
        \end{jieda}
        \item 已知$y=f(x)(x\zr)$是奇函数, 若当$x<0$时, $f(x)=x^{\frac{1}{3}}+2^x-1$, 则函数$f(x)$的解析式\\
        $f(x)$=\blkda[4]{$\begin{cases}
        x^{\frac{1}{3}}+2^x-1, x\leq 0\\
        x^{\frac{1}{3}}-2^{-x}+1, x>0
        \end{cases}$}. 
        \begin{jieda}
        当$x>0$时, $f(x)=-f(-x)=-((-x)^{\frac{1}{3}}+2^{-x}-1)=x^{\frac{1}{3}}-2^{-x}+1$. \\
        又上述两式当$x=0$时也满足$f(0)=0$. 
        \end{jieda}
    \end{enumerate}
    \item 方程的解和交点转化
    \begin{enumerate}
        \item 已知方程$|x| = ax+1$有一负解，且无正解，则实数$a$的取值范围是\blkda{$[1, +\infty)$ }
        \begin{jieda}
            将方程的根转化为$y = |x|$与$y = ax+1$的交点的横坐标.$y = ax+1$过$(0, 1)$，斜率为$a$，故当$a \in [1, +\infty)$时，两函数在$x > 0$时无交点，而在$x < 0$时有一个交点，满足要求.
        \end{jieda}
        \item 设$f(x)$是定义在$\bR$上的奇函数, 且对于任意的$x\zr$, $f(1+x)-f(1-x)=0$恒成立; 当$x\in[0,1]$时, $f(x)=2x$. 
        若方程$f(x)=ax$恰好有5个不同的解, 则实数$a$的取值范围为\blkda[2.5]{$(-\frac{2}{3},-\frac{2}{7})\cup\{\frac{2}{5}\}$}. 
        \begin{jieda}
        $f(x)$图像关于原点对称, 关于直线$x=1$对称, 可知周期$T=4$. \\
        $y = ax$的图像是过原点的直线，随着$a$的变化，其斜率变化. 绘制草图后可以判断要分为$a>0$与$a<0$两种情形, 
        不难求得$a=\frac{2}{5}$或$a\in(-\frac{2}{3},-\frac{2}{7})$. 
        \end{jieda}
        \item 已知$f(x) = x^2 + \frac{1}{x^2} + k (x + \frac{1}{x})+11, x \in (0, +\infty)$，若$f(x) = 0$有两个解，则实数$k$的取值范围是\blkda[4]{$(-\infty, -\frac{13}{2}) \cup \{-6\}$}
        \begin{jieda}
            首先考虑换元$t = x + \frac{1}{x}$，$f(x) = g(t) = t^2+k\cdot t + 9$，当$t = 2$时，只对应一个$x$，当$t > 2$时，对应着两个$x$. 故若$f(x)$在$(0, +\infty)$，有两个零点可以转化为$g(t)$在$(2, +\infty)$有一个零点，进一步转化为$y = -k$与$y = t + \frac{9}{t}$在$(2, +\infty)$有一个交点. 因此$k \in (-\infty, -\frac{13}{2}) \cup \{-6\}$ \\
            \ps \textbf{换元后及时判断新元的取值范围，并且分析清楚换元前后的对应关系！}
        \end{jieda}
        \item 若关于$x$的方程$\frac{|x|}{x+6} = kx^2$有四个不同的实数解，求实数$k$的取值范围.
        \begin{jieda}
            易知必有一根为0 \\
            参变分离有$\frac{1}{k} = \frac{(x+6)x^2}{|x|} = \left \{ \begin{array}{ll}
                \frac{(x+6)x^2}{x}, & x > 0 \\
                -\frac{(x+6)x^2}{x},  & x < 0
            \end{array}\right.$
            由于原方程必有一根为0，故可以转化为方程$\frac{1}{k} = \left \{ \begin{array}{ll}
                (x+6)x, & x > 0 \\
                -(x+6)x,  & x < 0
            \end{array}\right.$有三个不等实根.因此$\frac{1}{k} \in (0, 9)$，即$k \in (\frac{1}{9}, +\infty)$
        \end{jieda}
        \item 关于$x$的方程$k\cdot9^x-k\cdot3^{x+1}+6(k-5)=0$在$[0,2]$上有唯一解, 求实数$k$的取值范围. 
        \begin{jieda}
        令$t=3^x\in[1,9]$, $t$与$x$一一对应.于是关于$t$的方程$kt^2-3kt+6(k-5)=0$在$[1,9]$上有唯一解. \\
        于是$\frac{30}{k}=t^2-3t+6$. 结合图像可得$\frac{30}{k}\in(4,60]\cup\{\frac{15}{4}\}$, 即$k\in[\frac{1}{2},\frac{15}{2})\cup\{8\}$. \\
        \ps 此处参变分离的方式值得注意！
        \end{jieda}
        \item 已知关于$x$的方程$(2^x-1)^2-2\cdot\abs{2^x-1}+k=0$有三个根, 实数$k$的取值范围是\blkda{$(0, 1)$}. 
        \begin{jieda}
            换元$t = 2^x - 1$，$t \in (-1, +\infty)$，$t$与$x$一一对应. 故原方程有三个解即$t^2 - 2|t| + k = 0$有三个根，即$y = k$与$y = 2|t| - t^2$有三个交点.函数$y = 2|t| - t^2$在$(-1, 0)$单调递减，值域为$(0, 1)$，在$(0, 1)$单调递增，值域为$(0, 1)$，在$[1, +\infty)$单调递减，值域为$(-\infty, 1]$，故$k \in (0, 1)$时两函数有三个交点.
        \end{jieda}
    \end{enumerate}
    \item 利用必要条件简化问题
    \begin{enumerate}
        \item 已知函数$y=\frac{1}{4-2^x}$的图像关于点$P$对称, 则点$P$的坐标为\blkda{$(2,\frac{1}{8})$}
        \begin{jieda}
        设$P(a,b)$, 因为$x\neq 2$, 故$a=2$. \\
        $2b=f(1)+f(3)$, 于是$b=\frac{1}{8}$. \\
        \ps $f(x)$的对称中心是$(a, b)$或对称轴是$x = a$的必要条件均为$f(x)$的定义域关于$a$对称.
        \end{jieda}
        \item 已知二次函数 $f(x)=-\frac{1}{2} x^2+x$，是否存在实数 $m, n(m<n)$, 使得 $f(x)$ 的定义域为 $[m, n]$, 值域为 $[k m, k n](k>1)$.
        \begin{jieda}
        $f(x)=-\frac{1}{2} x^2+x=-\frac{1}{2} \times(x-1)^2+\frac{1}{2} \leq \frac{1}{2}$
        $\therefore k n \leq \frac{1}{2}, \therefore m<n<\frac{1}{2k}$
        又 $\because$ 函数的对称轴为 $x=1$
        $\therefore f(x)$ 在 $[m, n]$ 单调递增则有 $f(m)=k m, f(n)=k n$
        故$m, n$是$-\frac{1}{2}x^2 + x = kx$的两根，故$m = 2(1-k), n = 0$ \\
        \ps 根据值域判断函数在定义域内的单调性是简化，解决问题的突破口。\\
        常规解决方法是分析单调性，根据最值列方程组.
        \end{jieda}
        \item 若函数$f(x)$在区间$[a, b]$上的值域恰是$[\frac{1}{b}, \frac{1}{a}]$，则称区间$[a, b]$是$f(x)$的一个倒域区间. 求函数$g(x) = \left \{ \begin{array}{ll}
            -x^2+2x & 0 \leq x \leq 2 \\
            x^2 + 2x & -2 \leq x < 0
        \end{array}\right.$的所有倒域区间.
        \begin{jieda}
            $g(x)$的定义域为$[-2, 2]$，$g(x)$值域为$[-1, 1]$，故有$-2 \leq a < b \leq 2, -1 \leq \frac{1}{b} < \frac{1}{a} \leq 1$，解得$1 \leq a < b \leq 2$或$-2 \leq a < b \leq -1$.
            \begin{dingautolist}{172}
                \item  $1 \leq a < b \leq 2$，此时$g(x)$单调递减，有$\left \{ \begin{array}{l}
                     -a^2 + 2a = \frac{1}{a}  \\
                     -b^2 + 2b = \frac{1}{b} 
                \end{array} \right.$，即$a, b$是$-x^2 + 2x = \frac{1}{x}$的两根，即$a, b$是$-x^3 + 2x^2 -1 = 0$的两根，可以看出有一根为$1$，即$a, b$是$(x-1)(-x^2+x+1) = 0$的两根，故$a = 1, b = \frac{1+\sqrt{5}}{2}$
                \item $-2 \leq a < b \leq -1$，此种情况和前一种对称，可得$a = -\frac{1+\sqrt{5}}{2}, b = -1$
            \end{dingautolist}
            因此两个倒域区间为$[1, \frac{1+\sqrt{5}}{2}]$和$[-\frac{1+\sqrt{5}}{2}, -1]$
        \end{jieda}
        \item 若函数$f(x)$在区间$[m, n]$上的值域恰是$[m, n]$，则称区间$[m, n]$是$f(x)$的一个等域区间. 若区间$[m,n]$是函数$y = \frac{(a+1)x - \frac{1}{a}}{ax}(a \neq 0)$的一个等域区间.求$n-m$的最大值.
        \begin{jieda}
            $\frac{(a+1)x - \frac{1}{a}}{ax} = \frac{a+1}{a} - \frac{1}{a^2x}$，故函数为$y = \frac{a+1}{a} - \frac{1}{a^2x} (a \neq 0)$，易知函数的定义域为$(-\infty, 0) \cup (0, +\infty)$，故必有$m < n < 0$或$0 < m < n$. 故易知函数在$[m, n]$内严格单调递增，故$\left \{ \begin{array}{l}
                    \frac{a+1}{a} - \frac{1}{a^2m} = m \\
                    \frac{a+1}{a} - \frac{1}{a^2n} = n
                \end{array} \right.$，即$m, n$是$a^2x^2 - (a+1)ax + 1 = 0$的两根. 由韦达定理有$\left \{ \begin{array}{l}
                    m+n = \frac{a+1}{a}\\
                    m \cdot n = \frac{1}{a^2}
                \end{array} \right.$，由$\Delta > 0$得$((a+1)^2 - 4)a^2 > 0$，故$a \in (-\infty, -3]\cup [1, +\infty)$ \\
                $n - m = \sqrt{(m+n)^2 - 4mn} = \sqrt{(\frac{a+1}{a})^2 - \frac{4}{a^2}} = \sqrt{\frac{a^2+2a-3}{a^2}} \leq \frac{2\sqrt{3}}{3}$，当$a = 3$时取等号.
        \end{jieda}
    \end{enumerate}
    \item 反函数的定点问题
    \begin{enumerate}
        \item 设函数$f(x)$存在反函数$f^{-1}(x)$，且函数$y = x - f(x)$的图像过点$(1, 2)$，则函数$y = f^{-1}(x) - x$的图像一定过点\blkda{$(-1, 2)$}
        \begin{jieda}
            $y = x - f(x)$，代入$x = 1, y = 2$，则$f(1) = -1$，故$f^{-1}(-1) = 1$. $x = -1$代入$y = f^{-1}(x) - x$，$y = 2$，故过$(-1, 2)$ \\
            \ps 反函数过定点问题本质是$f(x)$过$(a, b)$，则反函数过$(b, a)$，在此基础上，通过代入特值解决问题。
        \end{jieda}
        \item 已知 $f(x)=\frac{a-x}{x-a-1}$, 若函数 $y=f(x)$ 的反函数图像的对称中心是 $(-1,3)$, 则实数 $a$ 的值是\blkda{$2$}.
        \begin{jieda}
            由反函数的对称中心知$f(x)$的对称中心为$(3, -1)$，因此$a+1 = 3$，故$a = 2$
        \end{jieda}
    \end{enumerate}
    \item 反函数平移问题
    \begin{enumerate}
        \item 如果函数$y = f(x), x \in \mathbf{R}$的反函数为$y = f^{-1}(x), x \in \mathbf{R}$，则函数$y = f^{-1}(x-1)(x\in \mathbf{R})$的反函数为\blkda[3]{$y = f(x)+1$}.
        \begin{jieda}
            反函数和原函数一个向下平移，令一个向左平移相同长度（长度为负，方向取反），两者可保持关于$y = x$对称.
        \end{jieda}
        \item 已知 $f(x)=\sqrt{x-1}$, 函数 $y=f(x)$ 的反函数为 $y=f^{-1}(x)$, 若 $f^{-1}(x)-a=f(x+a)$有实数根, 则实数 $a$ 的取值范围为\blkda{$[\frac{3}{4}, +\infty)$}.
        \begin{jieda}
            $f^{-1}(x) = x^2 + 1 (x \geq 0)$. \\
            设$g(x) = f(x+a)$，则$g^{-1}(x) = f^{-1}(x)-a$，因为$g(x)$与$g^{-1}(x)$均为增函数，因此其交点必然在$y = x$上。由图象可知，$a$的取值范围形式应为$[m, +\infty)$，$m$为边界情况对应的值.边界情况为$x^2 + 1 - a = x$只有一解，此时$a = \frac{3}{4}$，故$a \in [\frac{3}{4}, +\infty)$ 
        \end{jieda}
    \end{enumerate}
    \item 幂函数的性质
    \begin{enumerate}
        \item 如图是函数$y=x^{\frac{m}{n}}$($m,n\zn$且互质)的图像, 则\dotfill(\kh{C})
        \xx{$m,n$是奇数, 且$\frac{m}{n}<1$}
        {$m$是偶数, $n$是奇数, 且$\frac{m}{n}>1$}
        {$m$是偶数, $n$是奇数, 且$\frac{m}{n}<1$}
        {$m$是奇数, $n$是偶数, 且$\frac{m}{n}>1$}
        \begin{minipage}{\linewidth}
        \centering
        \includegraphics[totalheight=2cm]{./bodyxb/src/main42chap21mihanshu02.png}
        \end{minipage}
        \item 下列命题中, 真命题的是\dotfill(\kh{C})
        \xx{幂函数中不存在既不是奇函数, 又不是偶函数的函数;}
        {如果一个幂函数不是偶函数, 那么它一定是奇函数;}
        {图像不经过$(-1,1)$的幂函数, 一定不是偶函数;}
        {如果两个幂函数有三个公共点, 那么这两个函数一定相同.}
        \item 当$x \in (0, +\infty)$时，幂函数$f(x) = (m^2-m-1)x^{-5m-3}$为增函数，则实数$m =$\blkda{$-1$}
        \begin{jieda}
            由幂函数概念可知$m^2-m-1 = 1$，解得$m = -1 \mbox{或} 2$，又因为$f(x)$在$(0, +\infty)$为增函数，所以有$-5m-3 > 0$，因此$m = -1$
        \end{jieda}
    \end{enumerate}
    \item 指数函数和幂函数的交点
    \begin{enumerate}
        \item 函数$y = x^2$与$y = 2^x$的交点个数为\blkda{$3$}. 
        \begin{jieda}
            指数函数增长得很快（$a > 1$），当其从上方穿越了一个多项式函数后一定会从下方再次穿越这个多项式函数，$0 < a < 1$则是对称的情况.
        \end{jieda}
        \item 设常数$a \in \mathbf{R}$，则方程$|x+a| \cdot 2^x = 2018$的解的个数组成的集合是\blkda{$\{1,2,3\}$}
        \begin{jieda}
            方程变形为$|x+a| = 2018 \cdot 2^{-x}$，平移图像可知交点个数可能为1,2,3三种.
        \end{jieda}
        \item 设$min\{a, b, c\}$为$a, b, c$中最小值，则$y = min\{x+2, 10-x, 2^x\}$的最大值为\blkda{$6$}
        \begin{jieda}
            计算可知$y = 2^x$与$y = x+2$右侧交点为$(2, 4)$，故当$x > 2$时$x+2 < 2^x$，而$y = 10-x$与$y = 2-x$交于$(4, 6)$，结合图像可知$y = min\{x+2, 10-x, 2^x\}$的最大值为6.
        \end{jieda}
    \end{enumerate}
    \item 指数函数、对数函数翻折
    \begin{enumerate}
        \item 已知定义在$R^{+}$上的函数$f(x) = \left \{\begin{array}{ll}
        |log_3x - 1|, & 0 < x \leq 9 \\
        4-\sqrt{x}, & x > 9
        \end{array} \right.$，设$a, b, c$是三个互不相同的实数，满足$f(a) = f(b) = f(c)$，求$abc$的取值范围.
        \begin{jieda}
            不妨设$a < b < c$，易知$0 < a < 3 < b < 9< c$，因此有$1 - log_3a = log_3b - 1$，即$ab = 9$. \\
            故$f(a) \in (0, 1)$，因此$c \in (9, 16)$，因此有$abc \in (81, 144)$
        \end{jieda}
        \item 关于$x$的方程$|3^x-1| = 2a$只有一个解，则$a$的取值范围是\blkda[3]{$[\frac{1}{2}, +\infty) \cup \{0\}$}.
        \begin{jieda}
            画图即可，$2a = 1$时是函数的渐近线，可以取得. 不要忘记$a = 0$时候也是只有一个解！
        \end{jieda}
    \item 已知函数$f(x) = \left \{ \begin{array}{ll}
        |log_2x| & 0 < x < 2 \\
        x^2 - 8x + 13 & x \geq 2
    \end{array} \right.$，若$f(x) = a$有四个解$x_1, x_2, x_3, x_4$，则$x_1 + x_2 + x_3 + x_4$的取值范围是\blkda{$(10, \frac{21}{2})$}.
    \begin{jieda}
        $x_1 \cdot x_2 = 1, x_3 + x_4 = 8$，$x_1 + x_2 + x_3 + x_4 = 8 + x_2 + \frac{1}{x_2}, x_2 \in (1, 2)$，故$x_1 + x_2 + x_3 + x_4 \in (10, \frac{21}{2})$ \\
        \ps \textbf{对数函数翻折时，总有和常函数交点横坐标乘积为定值}.
    \end{jieda}
    \end{enumerate}
        \item 换元法求指数、对数型复合函数的值域、最值
        \begin{enumerate}
        \item 求函数$y=3\cdot4^x-20\cdot2^x+3\cdot4^{-x}-20\cdot2^{-x}$的最小值. 
        \begin{jieda}
        令$t=2^x+2^{-x}\in[2,+\infty)$, 则
        $y = 3t^2-20t-6 = 3(t-\frac{10}{3})^2-\frac{118}{3}$ \\
        故$y_{\min}=-\frac{118}{3}$, 当$t=\frac{10}{3}$即$2^x = \frac{1}{3}\mbox{或}3, x=log_23 \mbox{或}-log_23$时取到. 
        \end{jieda}
        \item 若$a^{2x}+\frac{1}{2}a^x-\frac{1}{2}\leq 0$($a>0$且$a\neq 1)$, 求函数$y=2a^{2x}-3a^x+4$的值域. 
        \begin{jieda}
        换元$t = a^x, t > 0$，原不等式化为$t^2+\frac{1}{2}t-\frac{1}{2}\leq 0$，解得$t\in(0,\frac{1}{2}]$. \\
        $y=2a^{2x}-3a^x+4=2(t-\frac{3}{4})^2+\frac{23}{8}\in[3,4)$. 
        \end{jieda}
        \item 已知$x$满足不等式$2(\log_{\frac{1}{2}}x)^2+7\log_{\frac{1}{2}}x+3\leq 0$, 
        求$f(x)=\log_2\frac{x}{2}\log_2\frac{x}{4}$的最大值与最小值.
        \begin{jieda}
        令$t=\log_2x$, 则$2t^2-7t+3\leq 0$, 解得$t\in[\frac{1}{2},3]$. \\
        故$y=(t-1)(t-2)=t^2-3t+2=(t-\frac{3}{2})^2-\frac{1}{4}\in[-\frac{1}{4},2]$. \\
        所以$f(x)_{\min}=f(2^{\frac{3}{2}})=-\frac{1}{4}$, $f(x)_{\max}=f(8)=2$. 
        \end{jieda}
    \end{enumerate}
    \item 对数运算 
    \begin{enumerate}
        \item 若$\log_427=m$, $\log_325=n$, 则$\lg2$可用$m,n$表示为\blkda{$\frac{3}{3+mn}$}. 
        \begin{jieda}
            $m = \frac{3log_23}{2}$，$n = \frac{2log_25}{log_23}$，$lg2 = \frac{1}{1+log_25}$，故$log_25 = \frac{m\cdot n}{2}$，$lg2 = \frac{3}{3+mn}$ \\
            \ps \textbf{对数运算表示类的题目，首先要统一底数，选取题干中的一个质因子作为底数，然后将所有的对数式都拆彻底（底数和真数互质），然后再凑.}
        \end{jieda}
        \item 求$\frac{1-2log_65 \cdot log_{10}3 \cdot log_{15}2}{log_65 \cdot log_{10}3 + log_{10}3 \cdot log_{15}2 + log_{15}2 \cdot log_65}$的值.
    \begin{jieda}
        考虑换底，$log_65 = \frac{log_25}{log_23+1}, log_{10}3 = \frac{log_23}{log_25+1}, log_{15}2 = \frac{1}{log_23 + log_25}$\\
        故设$a = log_23, b = log_25$\\
        原式等于$\frac{1 - 2 \cdot \frac{b}{a+1} \cdot \frac{a}{b+1} \cdot \frac{1}{a+b}}{\frac{b}{a+1}\cdot \frac{a}{b+1} + \frac{a}{b+1} \cdot \frac{1}{a+b} + \frac{1}{a+b} \cdot \frac{b}{a+1}} = \frac{(a+1)(b+1)(a+b) - 2ab}{ab(a+b) + a(a+1) + b(b+1)} = 1$
    \end{jieda}
    \end{enumerate}
    \item 自变量联动与不联动条件下的恒成立、有解问题
    \begin{enumerate}
        \item 已知函数$f(x) = x - 2$，$g(x) = mx^2-2mx+1(m \in \mathbf{R})$. 
        \begin{enumerate}[1.]
            \item 对任意的$x \in \mathbf{R}$，不等式$g(x) > f(x)$恒成立，求$m$取值范围.
            \item 对任意的$x_1 \in [1, 2], x_2 \in [-1, 0]$，使得$g(x_1) > f(x_2)$，求$m$取值范围.
            \item 对任意的$x_1 \in [1, 2]$，存在$x_2 \in [3, 4]$，使得$g(x_1) = f(x_2)$，求$m$取值范围.
            \item 对任意的$x_1 \in [1, 2]$，存在$x_2 \in [0, 2]$，使得$g(x_1) > f(x_2)$，求$m$取值范围.
        \end{enumerate}
        \begin{jieda}
            \begin{enumerate}[1.]
                \item 问题转化为$g(x) - f(x) > 0$在$\mathbf{R}$上恒成立，即$mx^2-(2m+1)x+3 > 0$恒成立. $m = 0$不满足要求，$m < 0$不满足要求，$m > 0$时应有$\Delta < 0$，即$m \in (\frac{2-\sqrt{3}}{2}, \frac{2+\sqrt{3}}{2})$
                \item 问题转化为$g(x)$在$[1, 2]$上的最小值大于$f(x)$在$[3, 4]$上的最大值. $f(x)$在$[-1, 0]$上的最大值为$-2$，故问题转化为$g(x) + 2 > 0$在$[1, 2]$上恒成立，设$h(x) = g(x) + 2 = mx^2-2mx + 3 = m(x-1)^2-(m-3)$，下面分析$h(x)$的最值
                \begin{dingautolist}{172}
                    \item $m > 0$，$h(x)$在$[1, 2]$内单调递增，因此最小值为$h(1) = 3-m$
                    \item $m = 0$，$h(x) = -1$，最小值为$-1$
                    \item $m < 0$，$h(x)$在$[1, 2]$内单调递减，因此最小值为$h(2) = -1$
                \end{dingautolist}
                故当$m \in (-\infty, 5)$时，满足要求.
                \item 问题转化为$g(x)$在$[1, 2]$上的值域是$f(x)$在$[3, 4]$上值域的子集. $f(x)$在$[3, 4]$上值域为$[1, 2]$，$g(x) = m(x-1)^2 + (1-m)$，下面分析$g(x)$的值域
                \begin{dingautolist}{172}
                    \item $m > 0$，$g(x)$在$[1, 2]$内单调递增，值域为$[1-m, 1]$
                    \item $m = 0$，$g(x) = 1$，值域为$\{1\}$
                    \item $m < 0$，$g(x)$在$[1, 2]$内单调递增，值域为$[1, 1-m]$
                \end{dingautolist}
                故当$m \in [-1, 0]$时满足要求.
                \item 问题转化为$g(x)$在$[1, 2]$上的最小值大于$f(x)$在$[0, 2]$上的最小值. $f(x)$在$[0, 2]$上最小值为-2，下面分析$g(x)$的最值
                \begin{dingautolist}{172}
                    \item $m > 0$，$g(x)$在$[1, 2]$内单调递增，因此最小值为$g(1) = 1-m$
                    \item $m = 0$，$g(x) = -1$，最小值为$-1$
                    \item $m < 0$，$g(x)$在$[1, 2]$内单调递减，因此最小值为$g(2) = 1$
                \end{dingautolist}
                故当$m \in (-\infty, 3)$时，满足要求.
            \end{enumerate}
            \ps \textbf{重要的是问题的转化！}
        \end{jieda}
    \end{enumerate}
    \item 同角正弦值、余弦值的基本关系
    \begin{enumerate}
        \item 已知$\dfrac{ cos\alpha}{1+ sin\alpha}=-\dfrac{1}{2}$, 则$\dfrac{ cos\alpha}{ sin\alpha-1}$=\blkda{2}. 
        \begin{jieda}
        $\dfrac{ cos\alpha}{1+ sin\alpha}= \dfrac{1-sin\alpha}{cos\alpha}$ \\
        \ps $\dfrac{ sin\alpha}{1+ cos\alpha}= \dfrac{1-cos\alpha}{sin\alpha}$ 
        \end{jieda}
        \item 已知 $ sin \alpha+ cos \alpha=\frac{1}{5}(0<\alpha<\pi)$, 求: \\
\begin{enumerate*}
    \item  $ sin \alpha- cos \alpha$;
    \item  $ sin ^3 \alpha+ cos ^3 \alpha$;
    \item $ sin ^4 \alpha+ cos ^4 \alpha$.
\end{enumerate*}
\begin{jieda}
    \begin{enumerate}
        \item $(sin\alpha+cos\alpha)^2 = 1 + 2sin\alpha cos\alpha = \frac{1}{25}$，因此$sin\alpha cos\alpha = -\frac{12}{25}$ \\
        因此有$(sin\alpha - cos\alpha)^2 = 1 - 2sin\alpha cos\alpha = \frac{49}{25}$ \\
        故$sin\alpha - cos\alpha = \pm \frac{7}{5}$
        \\
        根据$\alpha \in (0, \pi)$，故$sin \alpha - cos\alpha > -1$，因此$sin\alpha - cos\alpha = \frac{7}{5}$ \\
        \ps 可以根据$sin\alpha + cos\alpha$在不同象限的值域直接判断出来$\alpha$是第二象限角.
        \item $sin^3 \alpha + cos^3 \alpha = (sin\alpha + cos\alpha)(sin^2 \alpha + cos^2 \alpha - sin\alpha cos\alpha) = (\frac{1}{5})(1 + \frac{12}{25}) = \frac{37}{125}$ \\
        \ps 也可以在第一问的结果上构建方程组$\left \{ \begin{array}{}
            sin\alpha + cos\alpha = \frac{1}{5} \\
            sin\alpha - cos\alpha = \frac{7}{5} 
        \end{array}\right.$解得$\left \{ \begin{array}{}
            sin\alpha = \frac{4}{5} \\
            cos\alpha = -\frac{3}{5} 
        \end{array}\right.$后带入计算（下同）
        \item $ sin ^4 \alpha+ cos ^4 \alpha = ( sin ^2 \alpha+ cos ^2 \alpha)^2 - 2(sin \alpha cos\alpha)^2 = 1 - 2 \times (-\frac{12}{25})^2 = \frac{337}{625}$
    \end{enumerate}
\end{jieda}
    \end{enumerate}
    \item 切割化弦
    \begin{enumerate}
        \item 下列各式为正号的是\dotfill(\kh{C})
        \xx{$ cos2- sin2$}
        {$ cos2\cdot sin2$}
        {$ tan2\cdot\sec2$}
        {$ sin2\cdot tan2$}
        \begin{jieda}
        $2\in(\frac{\pi}{2},\pi)$, 故$ sin2>0$, $ cos2<0$, $ tan2<0$, $ cot2<0$, $\sec2<0$, $\csc2>0$. 
        \end{jieda}
        \item 若角$\alpha$是第三象限角, 则
        \ding{172} $ sin\alpha+ cos\alpha<0$; 
        \ding{173} $ tan\alpha- sin\alpha>0$; 
        \ding{174} $ cot\alpha\cdot\csc\alpha<0$; 
        \ding{175} $ sin\alpha\cdot\sec\alpha>0$. 
        其中正确的是\blkda{\ding{172}\ding{173}\ding{174}\ding{175}}.  
        
        \item 设$y= sin\alpha\cdot cos\alpha\cdot cot\alpha$, 且$\alpha$是象限角, 则$y$的符号为\dotfill(\kh{A})
        \xx{恒正}{恒负}{可能为0}{不确定}
        \item 已知$\theta\in\left(0,\frac{\pi}{2}\right)$, 且$\log_{ tan\theta+ cot\theta} sin\theta=-\frac{3}{4}$, 
        则$\log_{ tan\theta} cos\theta=$\blkda{$\frac{1}{2}$}.
        \begin{jieda}
        $( tan\theta+ cot\theta)^{-\frac{3}{4}} =(\frac{sin\theta}{cos\theta} + \frac{cos\theta}{sin\theta})^{-\frac{3}{4}} = (\frac{1}{sin\theta cos\theta})^{-\frac{3}{4}} = sin\theta$，即$sin\theta = cos^3\theta$. 故$ cos^2\theta= tan\theta$, 所以$\log_{ tan\theta} cos\theta=\frac{1}{2}$. 
        \end{jieda}
    \end{enumerate}
    \item 齐次同除
    \begin{enumerate}
        \item 已知$tan\alpha=\sqrt{2}$, 求下列各式的值: \\
        \begin{enumerate*}[1.]
        \item $\frac{cos\alpha-5sin\alpha}{3cos\alpha+sin\alpha}$; 
        \item $\frac{sin^2\alpha-sin\alpha cos\alpha-3cos^2\alpha}{5sin\alpha cos\alpha+sin^2\alpha+1}$; 
        \item $2sin^2\alpha-sin\alpha cos\alpha+cos^2\alpha$. 
        \end{enumerate*}
        \begin{jieda}
        \begin{enumerate}[1.]
        \item $\frac{cos\alpha-5sin\alpha}{3cos\alpha+sin\alpha} = \frac{1-5tan\alpha}{3+tan\alpha} = \frac{13-16\sqrt2}{7}$; 
        \item $\frac{sin^2\alpha-sin\alpha cos\alpha-3cos^2\alpha}{5sin\alpha cos\alpha+sin^2\alpha+1} = \frac{sin^2\alpha-sin\alpha cos\alpha-3cos^2\alpha}{5sin\alpha cos\alpha+2sin^2\alpha+cos^2\alpha} \\
        = \frac{tan^2\alpha-tan\alpha -3}{5 tan\alpha+2tan^2\alpha + 1} = -\frac{1}{5}$; 
        \item $2sin^2\alpha-sin\alpha cos\alpha+cos^2\alpha = \frac{2sin^2\alpha-sin\alpha cos\alpha+cos^2\alpha}{sin^2\alpha + cos^2\alpha} \\
        = \frac{2tan^2\alpha-tan\alpha+1}{tan^2\alpha + 1} = \frac{5-\sqrt{2}}{3}$. 
        \end{enumerate}
        \end{jieda}
    \end{enumerate}
    \item 利用诱导公式化简
    \begin{enumerate}
        \item 化简：$\frac{sin(2\pi - \alpha) cos(\pi + \alpha)cos(\frac{\pi}{2}+\alpha)cos(\frac{11\pi}{2}-\alpha)}{cos(\pi - \alpha)sin(3\pi-\alpha)sin(-\pi-\alpha)sin(\frac{9\pi}{2}+\alpha)}$
        \begin{jieda}
            原式 \begin{align*}
                & = \frac{sin(-\alpha) cos(\pi + \alpha)cos(\frac{\pi}{2}+\alpha)cos(-\frac{\pi}{2}-\alpha)}{cos(\pi - \alpha)sin(\pi-\alpha)sin(-\pi-\alpha)sin(\frac{\pi}{2}+\alpha)} \\
            & = \frac{(-sin\alpha) (-cos\alpha) (-sin\alpha)cos(\frac{\pi}{2} + \alpha)}{(-cos\alpha)(sin\alpha)(-sin(\pi+\alpha))(cos\alpha)} \\
            & = \frac{(-sin\alpha) (-cos\alpha) (-sin\alpha)(-sin\alpha)}{(-cos\alpha)(sin\alpha)(sin\alpha)(cos\alpha)} \\ & = -tan\alpha
            \end{align*}
            \ps 先“大化小”，找合适的终边相同的角，然后再“负化正”或直接应用公式.
        \end{jieda}
        \item \begin{enumerate}
            \item 已知$sin(\frac{\pi}{3}-\alpha) = \frac{1}{2}$，求$cos(\frac{\pi}{6} + \alpha)$的值.
            \item 已知$cos(\frac{5\pi}{12}+x) = \frac{1}{3}$，且$x \in (-\pi, -\frac{\pi}{2})$，求$cos(\frac{\pi}{12} - x)$的值.
        \end{enumerate}
        \begin{jieda}
            \begin{enumerate}
                \item 因为$(\frac{\pi}{3}-\alpha) + (\frac{\pi}{6} + \alpha) = \frac{\pi}{2}$，所以$cos(\frac{\pi}{6} + \alpha) = cos(\frac{\pi}{2} - (\frac{\pi}{3} - \alpha)) = sin(\frac{\pi}{3} - \alpha) = \frac{1}{2}$
                \item 因为$(\frac{5\pi}{12}+x) + (\frac{\pi}{12} - x) = \frac{\pi}{2}$，所以$cos(\frac{\pi}{12} - x) = cos(\frac{\pi}{2} - (\frac{5\pi}{12}+x)) = sin(\frac{5\pi}{12}+x)$ \\
                因为$x \in (-\pi, -\frac{\pi}{2})$，所以$\frac{5\pi}{12}+x \in (-\frac{7\pi}{12}, -\frac{\pi}{12})$，又因为$cos(\frac{5\pi}{12}+x) = \frac{1}{3}$，所以$\frac{5\pi}{12}+x$是第四象限角，故$sin(\frac{5\pi}{12}+x) = -\sqrt{1 - cos^2(\frac{5\pi}{12}+x)} = -\sqrt{1 - \frac{1}{9}} = \frac{-2\sqrt{2}}{3}$
            \end{enumerate}
            \ps 化简求值的问题，找到已知和未知之间特殊的数量关系可以简化问题.
        \end{jieda}
        \item 已知$sin\theta, cos\theta$是关于$x$的方程$x^2 - ax + a  = 0(a \in \mathbf{R})$的两根，
        \begin{enumerate*}
            \item 求$cos^3(\frac{\pi}{2} - \theta) + sin^3(\frac{\pi}{2} + \theta)$的值
            \item 求$tan(\pi - \theta) - \frac{1}{tan\theta}$的值
        \end{enumerate*}
        \begin{jieda}
            根据韦达定理有$\left \{ \begin{array}{l}
                 sin\theta + cos\theta = a  \\
                 sin\theta \cdot cos\theta = a
            \end{array} \right.$，由$\Delta \geq 0$有$a \in (-\infty, 0] \cup [4, +\infty)$ \\
            $(sin\theta + cos\theta)^2 = 1 + 2sin\theta \cdot cos\theta = 1 + 2a = a^2$，解得$a = 1 \pm \sqrt{2}$，结合$y = sinx + cosx$的取值范围以及$a \in (-\infty, 0] \cup [4, +\infty)$可知,$a = 1 - \sqrt{2}$
            \begin{enumerate}
            \item $cos^3(\frac{\pi}{2} - \theta) + sin^3(\frac{\pi}{2} + \theta) = cos^3\theta + sin^3\theta = (sin\theta + cos\theta)(sin^2\theta + cos^2\theta - sin\theta cos\theta) = a \cdot (1-a) = (1-\sqrt{2})(\sqrt{2}) = \sqrt{2} - 2$ 
            \item $tan(\pi - \theta) - \frac{1}{tan\theta} = -(tan\theta + \frac{1}{tan\theta}) = -(\frac{sin\theta}{cos\theta} + \frac{cos\theta}{sin\theta}) = - (\frac{1}{sin\theta \cdot cos\theta}) = -\frac{1}{a} = \frac{1}{\sqrt{2} - 1} = \sqrt{2}+1$
        \end{enumerate}
        \end{jieda}
        \item 在$\triangle ABC$中，若$sin(2\pi - A) = -\sqrt{2}sin(\pi - B)$，$\sqrt{3}cosA = -\sqrt{2}cos(\pi - B)$，求$\triangle ABC$ 的三个内角.
        \begin{jieda}
            $sin(2\pi - A) = -\sqrt{2}sin(\pi - B)$即$sinA =\sqrt{2}sinB$ \\
            $\sqrt{3}cosA = -\sqrt{2}cos(\pi - B)$即$\sqrt{3}cosA = \sqrt{2}cosB$ \\
            故有$sin^2A + cos^2B = 2sin^2B + \frac{2}{3}cos^2B = \frac{4}{3}sin^2B + \frac{2}{3} = 1$，即$sin^2B = \frac{1}{4}$ \\
            因此$sinB = \frac{1}{2}, sinA = \frac{\sqrt{2}}{2}$，因为$B$为三角形内角，$sinB = \frac{1}{2}$，所以$B = \frac{\pi}{6}$或$B = \frac{5\pi}{6}$，同理$A = \frac{\pi}{4}$或$\frac{3\pi}{4}$ \\
            若$B = \frac{5\pi}{6}$，则必有$A+B > pi$，不符合要求，因此$B = \frac{\pi}{6}, cosB = \frac{\sqrt{3}}{2}$，故$cosA = \frac{\sqrt{2}}{2}$，因此$A = \frac{\pi}{4}$. 故$C = \pi - A - B = \frac{7\pi}{12}$
        \end{jieda}
    \end{enumerate}
\end{enumerate}
% \chapter*{复习题} 
% \begin{enumerate}
%     \item 设集合$A = \{(x, y) | (x+3sin\alpha)^2 + (y+3cos\alpha)^2 = 1, \alpha \in \mathbf{R} \}$，$B = \{(x, y) | 3x+4y+10 = 0\}$，记$P = A \cap B$，则点集$P$的轨迹长度为\blkda{$4\sqrt{3}$}
%     \begin{jieda}
%         对于集合$A$考虑其共同特征的几何意义：点$(x, y)$到点$(-3sin\alpha, -3cos\alpha)$的距离为1。易知点$(-3sin\alpha, -3cos\alpha)$位于以原点为圆心，以3为半径的圆上，故$A$所对应的集合图像是以原点为圆心，以2为内径和4为外径的同心圆环及其内部。\\
%         而对于集合$B$，对应的是直线$3x+4y+10 = 0$，其到原点的距离为2。故其和A的小圆相切，和大圆相交，弦长为$4\sqrt{3}$，由此轨迹长度为$4\sqrt{3}$ \\
%     \end{jieda}
%     \item 已知$x, y \in \mathbf{R^+}$，则$\frac{xy+2yz}{x^2+y^2+z^2}$的最大值为\blkda{$\frac{\sqrt{5}}{2}$}
%     \begin{jieda}
%         分子的形式是乘积的和的形式，分母是平方和，易知分子分母之间可以用平均值不等式得到不等关系。观察可知分子中的$y$是复用的，因此考虑将分母中的$y^2$拆成两部分分别与$x^2$和$z^2$配对。因此设$a+b=1$，$\frac{xy+2yz}{x^2+ay^2+by^2+z^2}$，根据平均值不等式有$\frac{xy+2yz}{x^2+ay^2+by^2+z^2} \leq \frac{xy+2yz}{2\sqrt{a}xy + 2\sqrt{b}yz}$。\\
%         当$\sqrt{b} = 2\sqrt{a}$的时候，不等号右侧为常数。此时$a = 0.2, b = 0.8$，对应的常数为$\frac{\sqrt{5}}{2}$，因为平均值不等式是恒成立的不等关系，因此原式最大值为$\frac{\sqrt{5}}{2}$
%     \end{jieda}
%     \item 设$a, b, c$为正实数，求$\frac{a+3c}{a+2b+c} + \frac{4b}{a+b+2c} - \frac{8c}{a+b+3c}$的最小值\blkda[2.5]{$12\sqrt{2} - 17$}
%     \begin{jieda}
%         首先考虑通过换元将分母简化，设$x = a + 2b + c, y = a+b+2c, z = a+b+3c$。下一步通过$x, y, z$的线性组合来表示分子。易知$c = z-y$，故$x = a + 2b + (z-y), y = a + b + 2(z-y)$，即$a+2b = x+y-z, a+b=3y-2z$，故$b = x+z-2y$，故$a = 5y-x-3z$\\
%         由此原式变为$\frac{-x+2y}{x} + \frac{4(x+z-2y)}{y} - \frac{8(z-y)}{z}$，即$-1+2\frac{y}{x} + 4\frac{x}{y} + 4\frac{z}{y} -8 - 8 + 8\frac{y}{z}$，进一步整理得$-17 + \frac{2y}{x} + \frac{4x}{y} + \frac{4z}{y} + \frac{8y}{z}$ \\
%         根据平均值不等式有$-17 + \frac{2y}{x} + \frac{4x}{y} + \frac{4z}{y} + \frac{8y}{z} \geq -17 + 2\sqrt{8} + 2\sqrt{32} = 12\sqrt{2} - 17$，当$z = \sqrt{2}y, y = \sqrt{2}x$时取等号，令$x = 1, y = \sqrt{2}, z = 2$，此时$a = 5\sqrt{2}-7, b = 3-2\sqrt{2}, c = 2-\sqrt{2}$
%     \end{jieda}
%     \item 已知集合$A = \{a_1, a_2, a_3, \cdot \cdot \cdot, a_n \} ( n \geq 3, n \in \mathbf{N^*})$，集合$B =\{a_i + a_j | a_i \in A, a_j \in A, 1 \leq i < j \leq n \}$，则集合$B$中元素个数的最小值为\blkda[2]{$2n-3$}（用$n$的代数式表示）.
%     \begin{jieda}
%     不妨设$a_1 < a_2 < a_3 < ... < a_n$\\
%     此时必有$a_1+a_2 < a_1 + a_3 < ... < a_1 + a_n < a_2 + a_n < a_3 + a_n < ... < a_{n-1} + a_n$.\\
%     在这个不等关系中共有$2n-3$个元素，故$B$中元素个数$|B| \geq 2n-3$.\\
%     当$a_1, a_2, a_3, ..., a_n$构成公差非0的等差数列时，$a_m + a_n = a_p + a_q (m+n = p+q)$，必有$B$中元素个数为$2n-3$
%     \end{jieda}
%     \item 已知集合$S = \{1,2,3, ..., n\}$, 其中$n \in \mathbf{N^*}$, 集合$A,B$均为$S$的非空子集. 若集合$A$中的最大元素大于集合$B$中的最小元素, 则满足条件的有序集合对$(A, B)$共有\blkda[3]{$2^n \cdot (2^n-n-1)$}对.
%     \begin{jieda}
%     设集合$C$也是$S$的非空子集，且$A$中的最大元素小于等于集合$C$中的最小元素，则集合满足要求的$C$的个数与满足要求的集合$B$的个数相加为$S$的非空子集个数，即$2^n - 1$ \\
%     设$A$中最大元素为$a$，$C$的最小元素$c$，满足$c \geq a$ \\
%     故固定$a$时，满足要求的$C$的个数为$2^{n-a+1}-1$，因此满足要求的$B$的个数为$2^n - 2^{n-a+1}$，此时满足要求的$A$的个数为$2^{a-1}$，因此有序集合对$(A, B)$的个数为$(2^n - 2^{n-a+1}) \cdot 2^{a-1} = 2^{n+a-1} -2^n$ \\
%     遍历所有的$a$，对所有情况进行加和，得到$\sum_{a=1}^n (2^{n+a-1} -2^n) = \sum_{a=1}^n 2^{n+a-1} - n \cdot 2^n = 2^n \cdot (2^n-1) - n \cdot 2^n = 2^n \cdot (2^n-n-1)$
%     \end{jieda} 
%     \item 若函数$f(x) = \frac{(\sqrt{1008}x + \sqrt{1009})^2 + sin2018x}{2016x^2+2018}$的最大值为$M$，最小值为$m$，则$M + m = $\blkda{$1$}.
%     \begin{jieda}
%     首先考虑化简变形，令$a = \sqrt{1008}, b = \sqrt{1009}$，\\
%     $f(x) = \frac{(ax+b)^2 + sin(2b^2x)}{2(ax)^2+2b^2} = \frac{1}{2} + \frac{2abx + sin(2b^2x)}{2(ax)^2+2b^2}$ \\
%     $y = 2abx$为奇函数，$y = sin(2b^2x)$为奇函数，故$y = 2abx + sin(2b^2x)$为奇函数；$y = 2(ax)^2+2b^2$为偶函数；故$y = \frac{2abx + sin(2b^2x)}{2(ax)^2+2b^2}$为奇函数，故其最大值和最小值之和为0，因此$M + m = 1$
%     \end{jieda}
%     \item 已知函数$f(x) = x^3 + sinx (x \in \mathbf{R})$，函数$g(x)$满足$g(x)+g(2-x) = 0(x \in \mathbf{R})$，若函数$h(x) = f(x-1)-g(x)$恰有2019个零点，则所有这些零点之和为\blkda{$2019$}.
%     \begin{jieda}
%     易知$f(x)$为奇函数，故$f(x-1)$关于$(1, 0)$中心对称。又由$g(x)+g(2-x) = 0(x \in \mathbf{R})$可知$g(x)$关于$(1, 0)$中心对称。因此$h(x)$关于$(1, 0)$中心对称，因此其零点必然关于$1$对称（奇数个零点代表有一个零点为1），故所有零点之和为零点个数，2019.
%     \end{jieda}
%     \item 已知$f(x) = \frac{(2^x+1)^2}{2^x \cdot x}+1$在$[-2018, 0)\cup(0, 2018]$上最大值为$M$,最小值为$m$,则$M+m = $\blkda{$2$}.
%     \begin{jieda}
%     设$g(x) = \frac{(2^x+1)^2}{2^x \cdot x}$，则$g(x) = \frac{2^x+\frac{1}{2^x}+2}{x} = \frac{2^x+2^{-x}+2}{x}$，故$g(-x) = -g(x)$，即$g(x)$为奇函数，故$M+m = 2$
%     \end{jieda}
%     \item 若函数$y = f(x)$和$y = f(-x)$在$[m, n]$区间上的单调性相同，则把区间$[m, n]$叫做的“稳定区间”.已知区间$[1, 2023]$为函数$y = \abs{(\frac{1}{2})^x+a}$的“稳定区间”，则实数$a$的取值范围是\blkda{$[-2, -\frac{1}{2}]$}
%     \begin{jieda}
%         首先题干转化为$f(x)$在$[-n, -m]$和$[m, n]$两个区间上的单调性相反，由此易知$f(x)$在定义域内单调必然不会存在稳定区间。\\
%         设$g(x) = \abs{(\frac{1}{2})^x+a}$，若其存在稳定区间，则其必然在定义域内先减后增，即$a < 0$ \\
%         进一步，$g(x)$和$x$轴的交点应该在$[-m, m]$内才可以使得其具有稳定区间$[m, n]$，有$a \in [-2, -\frac{1}{2}]$
%     \end{jieda}
%     \item 设$F(x) = \frac{2^x}{a} - \frac{1}{2^x} + 2 - 2^{x-2} + \frac{a}{2^{x-2}}$，已知$F(x)$的最小值为$m$，且$m > 2 + \sqrt{7}$.求实数$a$的取值范围.
%     \begin{jieda}
%         $F(x) = (\frac{4-a}{4a})2^x + \frac{(4a-1)}{2^x} + 2$，易知$a$会影响$F(x)$的单调性
%         \begin{enumerate}
%             \item $a \geq 4$ \\
%             $\frac{4-a}{4a} \leq 0, (4a-1) > 0$，$F(x)$单调减，无最小值，不符合要求.
%             \item $\frac{1}{4} < a < 4$ \\
%             $\frac{4-a}{4a} > 0, (4a-1) > 0$，换元$t = 2^x (t > 0)$，则$F(x) = G(t) = (\frac{4-a}{4a})t + \frac{(4a-1)}{t} + 2$为关于$t$的对勾函数的一支.因此$G(t) \geq 2\sqrt{\frac{(4-a)(4a-1)}{4a}}+2$（$t = \sqrt{\frac{(4a-1)(4a)}{4-a}}$时取等号），\\
%             故$2\sqrt{\frac{(4-a)(4a-1)}{4a}} > \sqrt{7}$，解得$a \in (\frac{1}{2}, 2)$ 
%             \item $0 < a \leq \frac{1}{4}$ \\
%             $\frac{4-a}{4a} > 0, (4a-1) \leq 0$，$F(x)$单调增，无最小值，不符合要求.
%             \item $a = 0$\\
%             $F(x)$无意义.
%             \item $a < 0$ \\
%             $\frac{4-a}{4a} < 0, (4a-1) < 0$，这种情况和第二种情况对称，此时只有最大值无最小值.
%         \end{enumerate}
%     \end{jieda}
%     综上，$a \in (\frac{1}{2}, 2)$
%     \item 已知函数$f(x)$满足$f(x+1) = 1+\sqrt{2f(x)-f^2(x)}(x \in \mathbf{R})$，求$f(1) + f(2020)$的最大值.
%     \begin{jieda}
%         易知应分析函数的周期性 \\
%         首先分析$f(x)$取值范围，因为$f(x+1) = 1+\sqrt{2f(x)-f^2(x)}$,所以有$2f(x)-f^2(x) \geq 0$，故$f(x) \in [0, 2]$,$f(x+1) \in [1, 2]$，因此$f(x) \in [1, 2]$ \\
%         移项后两侧平方有$f^2(x+1) - 2f(x+1) + 1 = 2f(x) - f^2(x)$，若设$g(x) = 2f(x) - f^2(x)$，则$g(x) + g(x+1) = 1$。(利用结论，可以判断$g(x)$周期为2) \\
%         进一步迭代有$g(x+1) + g(x+2) = 1$，两式相减有$g(x)= g(x+2)$，故$g(x)$周期为2，即$2f(x)-f^2(x) = 2f(x+2)  - f^2(x+2)$ \\
%         移项有$2(f(x)-f(x+2))= (f(x) + f(x+2))(f(x)- f(x+2))$ \\
%         由$f(x) \in [1, 2]$知$f(x) + f(x+2) \geq 2$，故$f(x) = 1$或$f(x) = f(x+2)$ \\
%         故$f(1) + f(2020) = f(1) + f(0) = f(0) + 1 + \sqrt{2f(0)-f^2(0)}$ \\
%         设$t = f(0) (t \in [1, 2]), f(1) + f(2020) = h(t) = t + 1 + \sqrt{2t-t^2} = 2 + (t-1) + \sqrt{1-(t-1)^2}$ 
%         \begin{enumerate}
%             \item 方法1 \\
%             由柯西不等式知$2 + (t-1) + \sqrt{1-(t-1)^2}\leq 2 + 2\sqrt{(t-1)^2 + 2t-t^2} = 2+\sqrt{2}$，当$t = 1+\frac{\sqrt{2}}{2}$时取等.
%             \item 方法2 \\
%             设$k = (t-1)^2 (k \in [0, 1])$，$2 + (t-1) + \sqrt{1-(t-1)^2} = 2 + \sqrt{k} + \sqrt{1-k}$，设$y = \sqrt{k} + \sqrt{1-k}$，故$y^2 = 1 + 2\sqrt{k(1-k)} \leq 2$，故$2 + \sqrt{k} + \sqrt{1-k} \leq 2 + \sqrt{2}$，当$k = \frac{1}{2}$时取等.
%         \end{enumerate}
%         综上，$f(1) + f(2020) \leq 2+2\sqrt{2}$，当$f(0) = 1+\frac{\sqrt{2}}{2}$时取等.
         
%     \end{jieda}
%     \item 已知$m$是正整数，幂函数$f(x) = x^{1+m-m^2}$在$(0, +\infty)$上时严格增函数.
%     \begin{enumerate}
%         \item 求幂函数$f(x)$的解析式
%         \item 设$\varphi(x) = \frac{a[f(x)]^2+1}{bf(x)+c}$（$a, b, c$为整数），若$\varphi(x)$是奇函数，又$\varphi(1) = 2, \varphi(2) < 3$，求$a, b, c$的值.
%     \end{enumerate}
%     \begin{jieda}
%         \begin{enumerate}
%             \item 根据$f(x)$在$(0, +\infty)$上时严格增函数知$1+m-m^2>0$，又因为$m$是正整数，所以$m = 1$,故$f(x) = x$
%             \item $\varphi(x) = \frac{ax^2+1}{bx+c}$. 因为$\varphi(x)$是奇函数，故$0$不在定义域内，否则$\frac{1}{c}$必不为0，因此$c = 0$，经检验，此时$\varphi(x)$为奇函数. \\
%             $\varphi(1) = \frac{a+1}{b} = 2$，即$b = \frac{a+1}{2}$ \\
%             $\varphi(2) = \frac{4a+1}{2b} = \frac{4a+1}{a+1} = 4 - \frac{3}{a+1} < 3$，故$a \in \{0, 1\}$ \\
%             又因为$b$是整数，故$a = 1$\\
%             综上，$a = 1, b = 1, c = 0$
%         \end{enumerate}
%     \end{jieda}
%     \item 求函数$f(x) = x^2 + \sqrt{x^4 - 3x^2+2x+5}$的值域.
%     \begin{jieda}
%     $f(x) = x^2 + \sqrt{(x^2-2)^2+(x+1)^2}$，其几何意义为抛物线$y = x^2$上点$B$到$A(-1, 2)$的距离与$B$到$x$轴距离之和. 由图像可以判断出$f(x) \in [2, +\infty)$ 
%     \end{jieda}
%     \item 已知函数$f(x) = \left \{ \begin{array}{ll}
%         |log_2x| & 0 < x < 2 \\
%         x^2 - 8x + 13 & x \geq 2
%     \end{array} \right.$，若$f(x) = a$有四个解$x_1, x_2, x_3, x_4$，则$x_1 + x_2 + x_3 + x_4$的取值范围是\blkda{$(10, \frac{21}{2})$}.
%     \begin{jieda}
%         $x_1 \cdot x_2 = 1, x_3 + x_4 = 8$，$x_1 + x_2 + x_3 + x_4 = 8 + x_2 + \frac{1}{x_2}, x_2 \in (1, 2)$，故$x_1 + x_2 + x_3 + x_4 \in (10, \frac{21}{2})$
%     \end{jieda}
%     \item 已知$f(x) = log_2(4^x+2^{x+1}+1)-x$，若$f(2a-1)<f(a+3)$，则实数$a$的取值范围是\blkda{$(-\frac{2}{3}, 4)$}.
%     \begin{jieda}
%         $f(x) = log_2(\frac{4^x+2^{x+1}+1}{2^x}) = log_2(2^x+2^{-x}+2)$，为偶函数，“开口向上”，故$f(2a-1)<f(a+3)$转化为$|2a-1| < |a+3|$，解得$a \in (-\frac{2}{3}, 4)$.
%     \end{jieda}
% \end{enumerate}